\documentclass[draft,linenumbers]{agujournal2019}
\usepackage{url}
\usepackage{lineno}
\usepackage{soul}

\usepackage{xcolor}
\definecolor{default-linkcolor}{HTML}{A50000}
\definecolor{default-filecolor}{HTML}{A50000}
\definecolor{default-citecolor}{HTML}{4077C0}
\definecolor{default-urlcolor}{HTML}{4077C0}

\usepackage{hyperref}
\hypersetup{
  hidelinks,
  breaklinks=true
}


\providecommand{\tightlist}{\setlength{\itemsep}{0pt}\setlength{\parskip}{0pt}}
\providecommand{\listfigurename}{List of Figures}
\providecommand{\listtablename}{List of Tables}
\newcounter{none}
\setcounter{none}{0}

\makeatletter
\@ifundefined{pandocbounded}{
  \usepackage{graphicx}
  \newcommand{\pandocbounded}[1]{\setbox0=\hbox{#1}\ifdim\wd0>\linewidth\resizebox{\linewidth}{!}{#1}\else#1\fi}
}{}
\makeatother

\usepackage{longtable,booktabs,array}
\usepackage{calc}
\usepackage{etoolbox}
\makeatletter
\patchcmd\longtable{\par}{\if@noskipsec\mbox{}\fi\par}{}{}
\makeatother
\IfFileExists{footnotehyper.sty}{\usepackage{footnotehyper}}{\usepackage{footnote}}
\makesavenoteenv{longtable}

\NewDocumentCommand\citeproctext{}{}
\NewDocumentCommand\citeproc{mm}{\begingroup\def\citeproctext{#2}\cite{#1}\endgroup}
\makeatletter
 \let\@cite@ofmt\@firstofone
 \def\@biblabel#1{}
 \def\@cite#1#2{{#1\if@tempswa , #2\fi}}
\makeatother

\newlength{\cslhangindent}
\setlength{\cslhangindent}{1.5em}
\newlength{\csllabelwidth}
\setlength{\csllabelwidth}{3em}
\newenvironment{CSLReferences}[2]
 {\begin{list}{}{
  \setlength{\itemindent}{0pt}
  \setlength{\leftmargin}{0pt}
  \setlength{\parsep}{0pt}
  \ifodd #1
   \setlength{\leftmargin}{\cslhangindent}
   \setlength{\itemindent}{-\cslhangindent}
  \fi
  \setlength{\itemsep}{#2\baselineskip}
  }
  \item\relax}
 {\end{list}}
\usepackage{calc}
\newcommand{\CSLBlock}[1]{#1\hfill\break}
\newcommand{\CSLLeftMargin}[1]{\parbox[t]{\csllabelwidth}{#1}}
\newcommand{\CSLRightInline}[1]{\parbox[t]{\linewidth - \csllabelwidth}{#1}\break}
\newcommand{\CSLIndent}[1]{\hspace{\cslhangindent}#1}

\usepackage{amsmath}
\usepackage{amssymb}
\usepackage{amsfonts}
\usepackage{pdflscape}
\makeatletter
\@ifpackageloaded{subcaption}{}{\usepackage{subcaption}}
\@ifpackageloaded{caption}{}{\usepackage{caption}}
\captionsetup[subfigure]{margin=0.5em}
\AtBeginDocument{%
\renewcommand*\figurename{Figure}
\renewcommand*\tablename{Table}
}
\AtBeginDocument{%
\renewcommand*\listfigurename{List of Figures}
\renewcommand*\listtablename{List of Tables}
}
\newcounter{pandoccrossref@subfigures@footnote@counter}
\newenvironment{pandoccrossrefsubfigures}{%
\setcounter{pandoccrossref@subfigures@footnote@counter}{0}
\begin{figure}\centering%
\gdef\global@pandoccrossref@subfigures@footnotes{}%
\DeclareRobustCommand{\footnote}[1]{\footnotemark%
\stepcounter{pandoccrossref@subfigures@footnote@counter}%
\ifx\global@pandoccrossref@subfigures@footnotes\empty%
\gdef\global@pandoccrossref@subfigures@footnotes{{##1}}%
\else%
\g@addto@macro\global@pandoccrossref@subfigures@footnotes{, {##1}}%
\fi}}%
{\end{figure}%
\addtocounter{footnote}{-\value{pandoccrossref@subfigures@footnote@counter}}
\@for\f:=\global@pandoccrossref@subfigures@footnotes\do{\stepcounter{footnote}\footnotetext{\f}}%
\gdef\global@pandoccrossref@subfigures@footnotes{}}
\@ifpackageloaded{float}{}{\usepackage{float}}
\floatstyle{ruled}
\@ifundefined{c@chapter}{\newfloat{codelisting}{h}{lop}}{\newfloat{codelisting}{h}{lop}[chapter]}
\floatname{codelisting}{Listing}
\newcommand*\listoflistings{\listof{codelisting}{List of Listings}}
\makeatother

\journalname{JGR: Solid Earth}

\begin{document}


\title{Beyond Equilibrium II: Sluggish Ringwoodite Kinetics Amplify 660
km Discontinuity Depressions and Independently Promote Slab Stagnation}

\authors{Buchanan Kerswell\affil{1}, John Wheeler\affil{1}, Rene
Gassmöller\affil{2}, J. Huw Davies\affil{3}, Isabel
Papanagnou\affil{4}, Sanne Cottaar\affil{4}, David P. Dobson\affil{5}}

\affiliation{1}{Department of Earth, Oceans and Ecological Sciences,
University of Liverpool, Liverpool L69 3BX, UK}
\affiliation{2}{GEOMAR Helmholtz Centre for Ocean Research Kiel, Kiel,
Germany}
\affiliation{3}{School of Earth and Environmental Sciences, Cardiff
University, Park Place, Cardiff, CF10 3AT, UK}
\affiliation{4}{Bullard Laboratories, Department of Earth Sciences,
University of Cambridge, Cambridge CB3 0EZ, UK}
\affiliation{5}{Department of Earth Sciences, University College London,
London WC1E 6BS, UK}

\correspondingauthor{Buchanan Kerswell}{b.kerswell@liverpool.ac.uk}

\begin{keypoints}
\item Sluggish ringwoodite kinetics amplify thermal displacements 2--5
times, reconciling 30--60 km depressions with modest Clapeyron slopes.
\item Metastable ringwoodite drives progressive, persistent 660
stagnation, essential for interpreting diverse MTZ seismic and slab
structures.
\item Large experimental uncertainties and unresolved grain-size, water,
and Fe-Mg effects motivate new experiments on 660 reaction kinetics.
\end{keypoints}

\begin{abstract}
The 660 km seismic discontinuity (`660') beneath some subduction zones
exhibits 30--60 km deep, broad transitions far exceeding the \(\sim\)
3--16 km displacements predicted by cold slab temperatures and
experimentally-determined Clapeyron slopes (-0.4 to -1.3 MPa
K\(^{-1}\)). Recent studies emphasize compositional or hydration
effects, leaving the impact of diffusion-controlled ringwoodite
decomposition kinetics unexplored. Here we couple experimentally
calibrated olivine \(\Leftrightarrow\) wadsleyite \(\Leftrightarrow\)
ringwoodite \(\Leftrightarrow\) post-spinel kinetics to 2D compressible
subduction simulations, varying the ringwoodite \(\Leftrightarrow\)
post-spinel kinetic prefactor \(Z_\mathrm{ri}\) alongside slab
temperature, strength, inflow velocity, and lower-mantle viscosity jump.
Under the most sluggish kinetics permitted by experimental uncertainty,
delayed ringwoodite decomposition amplifies thermal displacements of the
660 by 2--5 times and independently promotes slab stagnation, even in
warm slabs. As reaction rates decrease, buoyant metastable ringwoodite
drives slabs from penetration to stagnation, slowing descent up to
tenfold, broadening the 660 to 34 km, and adding up to 64 km of kinetic
displacement (80 km total). Kinetic stagnation at the 660 is progressive
and persistent, a behavior essential for interpreting transition-zone
seismic signatures and slab morphologies. Under nominal ringwoodite
kinetics (\(Z_\mathrm{ri}\) best-fit to experimental data), rapid
ringwoodite decomposition holds the 660 near its equilibrium depth
(6--12 km kinetic displacement), even in cold slabs. While nominal
ringwoodite kinetics imply mostly thermal and/or compositional effects
displacing the 660, sluggish ringwoodite kinetics reconcile 30--60 km
depressions with modest Clapeyron slopes, motivating new experiments to
address large laboratory uncertainties and natural variations in
grain-size, water, and Fe-Mg content in slabs.
\end{abstract}

\section*{Plain Language Summary}
Near 660 kilometers beneath Earth's surface, the mineral ringwoodite
breaks down into denser lower-mantle minerals, creating a global seismic
boundary. Beneath some sinking tectonic plates (slabs), seismic studies
image this boundary 30 to 60 kilometers deeper and much thicker than
normal, far exceeding the 3 to 16 kilometer shift expected from cold
slab temperatures alone. Recent explanations emphasize rock composition
or water, but here we use computer simulations of sinking slabs to test
the impact of ringwoodite reaction speeds (kinetics). Under the slowest
reaction speeds permitted by experimental uncertainty, delayed
ringwoodite breakdown amplifies thermal displacements of the 660
boundary by 2--5 times and independently promotes stalling, even in warm
slabs. As reaction speeds decrease, buoyant untransformed ringwoodite
slows descent up to tenfold, progressively driving slabs from sinking to
persistent stalling, broadening the boundary to 34 kilometers, and
adding up to 64 kilometers of delay (80 kilometers total). Under nominal
reaction speeds (best-fit to experimental data), rapid ringwoodite
breakdown holds the seismic boundary near its equilibrium depth, even in
cold slabs. While nominal reaction speeds imply mostly thermal or
compositional shifts, sluggish kinetics independently explain 30 to 60
kilometer displacements, motivating new experiments to reduce laboratory
uncertainties.

\clearpage

\section*{Definition of Symbols}\label{sec:definition-of-symbols}
\addcontentsline{toc}{section}{Definition of Symbols}

{\def\LTcaptype{none} % do not increment counter
\begin{longtable}[]{@{}
  >{\raggedright\arraybackslash}p{(\linewidth - 6\tabcolsep) * \real{0.3214}}
  >{\raggedright\arraybackslash}p{(\linewidth - 6\tabcolsep) * \real{0.2143}}
  >{\raggedright\arraybackslash}p{(\linewidth - 6\tabcolsep) * \real{0.1429}}
  >{\raggedright\arraybackslash}p{(\linewidth - 6\tabcolsep) * \real{0.3214}}@{}}
\toprule\noalign{}
\begin{minipage}[b]{\linewidth}\raggedright
Parameter
\end{minipage} & \begin{minipage}[b]{\linewidth}\raggedright
Symbol
\end{minipage} & \begin{minipage}[b]{\linewidth}\raggedright
Unit
\end{minipage} & \begin{minipage}[b]{\linewidth}\raggedright
Equations
\end{minipage} \\
\midrule\noalign{}
\endhead
\bottomrule\noalign{}
\endlastfoot
Activation energy & \(E^{\ast}\) & J mol\(^{-1}\) &
\ref{eq:growth-rate-660}, \ref{eq:reaction-rate-660} \\
Activation enthalpy & \(H^{\ast}\) & J mol\(^{-1}\) &
\ref{eq:reaction-rate-410} \\
Activation volume & \(V^{\ast}\) & m\(^3\) mol\(^{-1}\) &
\ref{eq:reaction-rate-410} \\
Activation factor (rheology) & \(B\) & - & \ref{eq:rheological-model} \\
Clapeyron slope & \(\gamma\) & Pa K\(^{-1}\) &
\ref{eq:thermal-displacement} \\
Compressibility (reference) & \(\bar{\beta}\) & Pa\(^{-1}\) &
\ref{eq:density-ala} \\
Density & \(\rho\) & kg m\(^{-3}\) &
\ref{eq:navier-stokes-no-inertia}, \ref{eq:continuity-compressible}, \ref{eq:energy}, \ref{eq:continuity-expanded}, \ref{eq:density-ala}, \ref{eq:thermal-displacement} \\
Density (reference) & \(\bar{\rho}\) & kg m\(^{-3}\) &
\ref{eq:adiabatic-pressure}, \ref{eq:density-ala} \\
Density (dynamic) & \(\hat{\rho}\) & kg m\(^{-3}\) & - \\
Depth coordinate & \(y\) & m &
\ref{eq:adiabatic-temperature}, \ref{eq:adiabatic-pressure} \\
Deviatoric stress tensor & \(\sigma^{\prime}\) & Pa &
\ref{eq:navier-stokes-no-inertia}, \ref{eq:energy} \\
Deviatoric strain rate tensor & \(\dot{\epsilon}^{\prime}\) & s\(^{-1}\)
& \ref{eq:energy} \\
Displacement (total) & \(\delta z\) & km &
\ref{eq:kinetic-displacement} \\
Displacement (kinetic contribution) & \(\delta z_k\) & km &
\ref{eq:kinetic-displacement} \\
Displacement (thermal contribution) & \(\delta z_t\) & km &
\ref{eq:thermal-displacement}, \ref{eq:kinetic-displacement} \\
Gas constant & \(R\) & J mol\(^{-1}\) K\(^{-1}\) &
\ref{eq:reaction-rate-410}, \ref{eq:growth-rate-660}, \ref{eq:reaction-rate-660} \\
Gibbs width factor & \(w\) & J mol\(^{-1}\) &
\ref{eq:transition-function} \\
Grain size & \(d\) & m & - \\
Gravitational acceleration & \(g\) & m s\(^{-2}\) &
\ref{eq:navier-stokes-no-inertia}, \ref{eq:adiabatic-temperature}, \ref{eq:adiabatic-pressure}, \ref{eq:thermal-displacement} \\
Growth rate & \(\dot{x}\) & m s\(^{-1}\) &
\ref{eq:volume-fraction}, \ref{eq:volume-transformation-rate}, \ref{eq:growth-rate-660} \\
Growth rate prefactor & \(K^{\prime}\) & m\(^7\) J\(^{-2}\) s\(^{-1}\) &
\ref{eq:growth-rate-660} \\
Interlamellar spacing & \(\lambda\) & m & - \\
Kinetic prefactor (olivine \(\Leftrightarrow\) wadsleyite) &
\(Z_\mathrm{ol}\) & K s\(^{-1}\) & \ref{eq:reaction-rate-410} \\
Kinetic prefactor (wadsleyite \(\Leftrightarrow\) ringwoodite) &
\(Z_\mathrm{wd}\) & K s\(^{-1}\) & \ref{eq:reaction-rate-410} \\
Kinetic prefactor (ringwoodite \(\Leftrightarrow\) bridgmanite +
periclase) & \(Z_\mathrm{ri}\) & mol\(^2\) J\(^{-2}\) s\(^{-1}\) &
\ref{eq:reaction-rate-660} \\
Latent heat & \(Q_L\) & J kg\(^{-1}\) & \ref{eq:energy} \\
Molar entropy & \(\bar{S}\) & J mol\(^{-1}\) K\(^{-1}\) &
\ref{eq:gibbs} \\
Molar Gibbs free energy & \(\Delta G\); \(\bar{G}\) & J mol\(^{-1}\) &
\ref{eq:reaction-rate-410}, \ref{eq:gibbs}, \ref{eq:growth-rate-660}, \ref{eq:reaction-rate-660}, \ref{eq:transition-function} \\
Molar volume & \(\bar{V}\) & m\(^{3}\) mol\(^{-1}\) &
\ref{eq:gibbs}, \ref{eq:growth-rate-660} \\
Nucleation site factor & \(N\) & m\(^{-1}\) &
\ref{eq:volume-fraction}, \ref{eq:volume-transformation-rate} \\
Pressure & \(P\) & Pa &
\ref{eq:navier-stokes-no-inertia}, \ref{eq:energy}, \ref{eq:reaction-rate-410} \\
Pressure (reference) & \(\bar{P}\) & Pa & \ref{eq:adiabatic-pressure} \\
Pressure (dynamic) & \(\hat{P}\) & Pa &
\ref{eq:density-ala}, \ref{eq:gibbs} \\
Reaction rate & \(\frac{dX}{dt}\); \(\dot{X}\) & s\(^{-1}\) &
\ref{eq:volume-transformation-rate}, \ref{eq:reaction-rate-410}, \ref{eq:reaction-rate-660}, \ref{eq:composition} \\
Reaction rate (maximum at 660) & \(\dot{X}_\mathrm{max}\) & Ma\(^{-1}\)
& \ref{eq:reaction-timescale}, \ref{eq:characteristic-length-scale} \\
Reaction completion timescale & \(\tau\) & Ma &
\ref{eq:reaction-timescale} \\
Reaction completion length scale & \(L\) & km &
\ref{eq:characteristic-length-scale} \\
Specific heat capacity (reference) & \(\bar{C}_p\) & J kg\(^{-1}\)
K\(^{-1}\) & \ref{eq:energy}, \ref{eq:adiabatic-temperature} \\
Temperature & \(T\) & K &
\ref{eq:energy}, \ref{eq:reaction-rate-410}, \ref{eq:growth-rate-660}, \ref{eq:reaction-rate-660} \\
Temperature (reference) & \(\bar{T}\) & K &
\ref{eq:adiabatic-temperature}, \ref{eq:rheological-model} \\
Temperature (dynamic) & \(\hat{T}\) & K &
\ref{eq:density-ala}, \ref{eq:gibbs}, \ref{eq:rheological-model}, \ref{eq:thermal-displacement} \\
Temperature anomaly (inflow) & \(\Delta T_0\) & K & - \\
Thermal conductivity (reference) & \(\bar{k}\) & W m\(^{-1}\) K\(^{-1}\)
& \ref{eq:energy} \\
Thermal expansivity (reference) & \(\bar{\alpha}\) & K\(^{-1}\) &
\ref{eq:energy}, \ref{eq:adiabatic-temperature}, \ref{eq:density-ala} \\
Time & \(t\) & s &
\ref{eq:continuity-compressible}, \ref{eq:energy}, \ref{eq:continuity-expanded}, \ref{eq:volume-fraction}, \ref{eq:composition} \\
Transition function & \(f_j\) & - &
\ref{eq:transition-function}, \ref{eq:equilibrium-fractions} \\
Velocity (field) & \(\vec{u}\) & m s\(^{-1}\) &
\ref{eq:continuity-compressible}, \ref{eq:energy}, \ref{eq:continuity-expanded}, \ref{eq:composition} \\
Velocity (boundary) & \(\vec{u}_\mathrm{in}\) & cm yr\(^{-1}\) & - \\
Velocity (compressional wave) & \(V_p\) & km s\(^{-1}\) & - \\
Velocity (maximum through 660) & \(\vec{v}_\mathrm{max}\) & cm
yr\(^{-1}\) & \ref{eq:characteristic-length-scale} \\
Velocity (vertical component) & \(\vec{u}_y\) & cm yr\(^{-1}\) & - \\
Viscosity & \(\eta\) & Pa s & \ref{eq:rheological-model} \\
Viscosity (reference) & \(\bar{\eta}\) & Pa s &
\ref{eq:rheological-model} \\
Viscosity jump factor (lower mantle) & \(b\) & - &
\ref{eq:rheological-model}, \ref{eq:rheological-jump-factor} \\
Viscosity prefactor (phase \(i\)) & \(\nu\) & - &
\ref{eq:rheological-jump-factor} \\
Volume fraction (actual) & \(X\) & - &
\ref{eq:volume-fraction}, \ref{eq:volume-transformation-rate}, \ref{eq:reaction-rate-410}, \ref{eq:reaction-rate-660}, \ref{eq:composition} \\
Volume fraction (equilibrium) & \(\phi\) & - &
\ref{eq:equilibrium-fractions}, \ref{eq:rheological-jump-factor} \\
Width (660 transition) & \(\delta w\) & km & - \\
\end{longtable}
}

\clearpage

\section{Introduction}\label{sec:introduction}

Earth's 660 km seismic discontinuity (hereafter the `660') marks the
base of the mantle transition zone (MTZ), where ringwoodite decomposes
into bridgmanite and periclase (the `post-spinel' reaction). Two
shallower discontinuities, near 410 and 520 km depth (hereafter the
`410' and `520'), mark the olivine \(\Leftrightarrow\) wadsleyite and
wadsleyite \(\Leftrightarrow\) ringwoodite transitions that further
define the MTZ. High-pressure experiments constrain the Clapeyron slope
of the post-spinel reaction to -0.4 to -1.3 MPa K\(^{-1}\) (Fei et al.,
2004; Ishii et al., 2018; {Katsura et al.}, 2003; Kojitani et al., 2016;
Litasov et al., 2005), while steeper values up to -3 MPa K\(^{-1}\) were
obtained from earlier calibrations (Bina \& Helffrich, 1994; Ito \&
Takahashi, 1989) or hydrous compositions (Ghosh et al., 2013). This
negative Clapeyron slope deepens the phase transition beneath cold
subducting slabs and has long been recognized as a potential barrier to
whole-mantle convection (Ita \& Stixrude, 1992; Ringwood, 1991).

Equilibrium thermodynamics explains the 660's first-order behavior, but
two observations show it is incomplete. First, 660 structure varies more
than thermal effects alone predict. While global seismic studies show
patterns broadly consistent with the thermal modulation expected from a
negative Clapeyron slope (e.g., Flanagan \& Shearer, 1998; Lebedev et
al., 2002), amplitude variations point to additional compositional
heterogeneity ({Goes et al.}, 2022; Waszek et al., 2021; {Yu et al.},
2023), and local 660 structures in many subduction zones show
displacements of up to 30--60 km (e.g., Fukao \& Obayashi, 2013; Ba et
al., 2025; Cottaar \& Deuss, 2016; Tang et al., 2014) that far exceed
the \(\sim\) 3--16 km thermal displacement predicted by modest Clapeyron
slopes of -0.4 to -1.3 MPa K\(^{-1}\). Proposed explanations include
water in ringwoodite (Muir et al., 2021), the akimotoite
\(\Leftrightarrow\) bridgmanite transition ({Chanyshev et al.}, 2022;
Cottaar \& Deuss, 2016), and grain-size-dependent viscosity reduction
(Kubo et al., 2000; Mao \& Zhong, 2021; Yamazaki et al., 2005). Second,
slab dynamics at the 660 depend on more than the Clapeyron slope alone.
Numerical geodynamic studies show that slab stagnation above the 660
requires a negative Clapeyron slope, a lower-mantle viscosity increase,
and trench retreat acting jointly ({Agrusta et al.}, 2017; Garel et al.,
2014; {Goes et al.}, 2017). Whether reaction kinetics can independently
explain these observations has not been evaluated. Few studies have
moved beyond thermodynamic equilibrium to test how much microscale
reaction kinetics influence macroscale flow dynamics and seismic
structures across the MTZ. Building from previous work on metastable
olivine (Däßler et al., 1996; Devaux et al., 1997; Hosoya et al., 2005;
Kerswell et al., 2026; Li \& Billen, 2026; Riedel \& Karato, 1997; Rubie
\& Ross, 1994), the reaction-kinetics hypothesis offers a critical
avenue of inquiry that may independently account for both deep, broad
660 discontinuities and slab stagnation where equilibrium models fall
short.

Laboratory experiments establish that the 410 olivine
\(\Leftrightarrow\) wadsleyite reaction and the 660 post-spinel reaction
are governed by fundamentally different physical mechanisms (Dobson \&
Mariani, 2014; Hosoya et al., 2005; Kubo et al., 2002; Lessing et al.,
2022), which translate into different semi-empirical kinetic rate
equations with distinct mathematical formulations and parameter
dependencies. In a companion study, Kerswell et al. (2026) coupled
\emph{interface-controlled} olivine \(\Leftrightarrow\) wadsleyite
kinetics to compressible mantle flow simulations and identified three
distinct kinetic regimes in cold slabs with diagnostic seismic
signatures at the 410. Because \emph{diffusion-controlled} ringwoodite
decomposition kinetics take a different mathematical form, it is
expected to produce qualitatively different dynamic behavior at the 660.
The reverse post-spinel reaction in upwelling mantle follows a distinct
mechanism and proceeds much more slowly below 2000 K (Dobson \& Mariani,
2014; Lessing et al., 2022), so we restrict this study to forward
ringwoodite decomposition in subducting slabs. Whether
interface-controlled and diffusion-controlled mechanisms produce
different flow dynamics and discontinuity structures in subducting
slabs, and under what thermal and kinetic conditions reaction kinetics
at the 660 control large-scale mantle flow, has not been systematically
evaluated in geodynamic simulations of Earth's upper mantle.

Here, we couple experimentally calibrated olivine \(\Leftrightarrow\)
wadsleyite \(\Leftrightarrow\) ringwoodite \(\Leftrightarrow\)
post-spinel kinetics (after Hosoya et al., 2005; Kubo et al., 2002) to
compressible mantle flow simulations using ASPECT (Heister et al., 2017;
Kronbichler et al., 2012). Our simulations simultaneously resolve the
410, 520, and 660 transitions, enabling the first direct comparison
between reaction mechanisms in a consistent dynamic framework. We
hypothesize that ringwoodite decomposition kinetics control whether
slabs penetrate or stagnate at the 660, depending on how fast reaction
rates are relative to slab descent. Sluggish ringwoodite decomposition
should sustain a buoyant metastable zone that impedes descent, whereas
faster reaction rates should facilitate continuous penetration. We
further hypothesize that sluggish ringwoodite decomposition should
thicken and flatten slabs above the 660, amplifying thermal
displacements into seismically broad, deeply depressed discontinuities.
These hypotheses motivate three specific questions:

\begin{enumerate}
\def\labelenumi{\arabic{enumi}.}
\tightlist
\item
  Under what thermal and kinetic conditions do ringwoodite decomposition
  kinetics control 660 structure and slab dynamics, and how do these
  effects differ from those at the 410?
\item
  How do rheological strength contrasts and slab kinematics modulate
  thermal and kinetic contributions to 660 structure?
\item
  Can seismic observations of 660 structure independently constrain
  ringwoodite decomposition kinetics?
\end{enumerate}

To separate kinetic inhibition from thermal displacement, we compare two
end-member simulation suites that both fall within the 95\% confidence
envelope of the Kubo et al. (2002) calibration. We define a
\emph{Sluggish Kinetics / Warm Slabs} suite (SW) with 500 K slab thermal
anomalies and sluggish kinetic prefactors matching the effective dry,
coarse-grained conditions of Kerswell et al. (2026). We also define a
\emph{Nominal Kinetics / Cold Slabs} suite (NC) with 1000 K slab thermal
anomalies and nominal kinetic prefactors (\(Z_\mathrm{ri}\) best-fit to
experimental data), analogous to classic studies of metastable olivine
(Däßler et al., 1996; Devaux et al., 1997; Hosoya et al., 2005; Riedel
\& Karato, 1997; Rubie \& Ross, 1994). Across both suites, we vary the
ringwoodite \(\Leftrightarrow\) post-spinel kinetic prefactor
\(Z_\mathrm{ri}\) over six orders of magnitude alongside rheological
strength contrasts, inflow velocities, and lower-mantle viscosity
contrasts, and decompose total 660 displacement into its thermal and
kinetic contributions. In the NC suite, the 660 stays near equilibrium
across all six orders of magnitude in \(Z_\mathrm{ri}\). As
\(Z_\mathrm{ri}\) decreases within the SW suite, slabs progressively
slow and amplify the thermal 660 signal, without any abrupt thresholds
like the 410 (Kerswell et al., 2026). We discuss why the two reaction
mechanisms respond differently to cold slab temperatures, how jointly
observing the 410 and 660 in the same slab column helps separate kinetic
from thermal and compositional contributions to MTZ structure, and why
new experiments that reduce laboratory uncertainties are necessary to
resolve the role of kinetics in Earth's mantle.

\section{Methods}\label{sec:methods}

\subsection{Governing Equations for Compressible Mantle
Flow}\label{sec:governing-equations-for-compressible-mantle-flow}

We solve the same governing equations for compressible mantle flow as
Kerswell et al. (2026), summarized here for completeness. We simulate
mantle flow using the finite-element geodynamic code ASPECT v3.1.0
(Bangerth et al., 2018, 2026; Clevenger \& Heister, 2021; Gassmöller et
al., 2018) to determine the velocity \(\vec{u}\), pressure \(P\), and
temperature \(T\) fields that satisfy the following equations:

\begin{equation}\protect\phantomsection\label{eq:navier-stokes-no-inertia}{\nabla P - \nabla \cdot \sigma^{\prime} = \rho\, g}\end{equation}

\begin{equation}\protect\phantomsection\label{eq:continuity-compressible}{\frac{\partial \rho}{\partial t} + \nabla \cdot (\rho\, \vec{u}) = 0}\end{equation}

\begin{equation}\protect\phantomsection\label{eq:energy}{\rho\, \bar{C}_p \left(\frac{\partial T}{\partial t} + \vec{u} \cdot \nabla T \right) - \nabla \cdot \left(\bar{k}\, \nabla T \right) = \sigma^{\prime} : \dot{\epsilon}^{\prime} + \bar{\alpha}\, T \left(\vec{u} \cdot \nabla P \right) + Q_L}\end{equation}

where \(\sigma^{\prime}\) is the deviatoric stress tensor,
\(\dot{\epsilon}^{\prime}\) is the deviatoric strain rate tensor,
\(\rho\) is density, \(g\) is gravitational acceleration, \(t\) is time,
and \(\bar{C}_p\), \(\bar{k}\), \(\bar{\alpha}\) are reference specific
heat capacity, thermal conductivity, and thermal expansivity,
respectively. The term \(Q_L\) captures latent heat from phase
transitions. We adopt the projected density approximation (Gassmöller et
al., 2020), reformulating Equation~\ref{eq:continuity-compressible} as:

\begin{equation}\protect\phantomsection\label{eq:continuity-expanded}{\frac{1}{\rho} \frac{\partial \rho}{\partial t} + \nabla \cdot \vec{u} + \left(\frac{1}{\rho} \nabla \rho \right) \cdot \vec{u} = 0}\end{equation}

The projected density \(\rho(T, P, X)\) varies with temperature,
pressure, and volume fractions of the relevant mineral phases, ensuring
that local density changes from both pressure-temperature (PT)
variations and phase transitions influence the flow through buoyancy,
volumetric expansion, and compression. See Gassmöller et al. (2020) for
a detailed discussion of this compressible formulation.

\subsection{Numerical Setup}\label{sec:numerical-setup}

\subsubsection{Adiabatic Reference
Conditions}\label{sec:adiabatic-reference-conditions}

We compute adiabatic reference conditions using the same approach as
Kerswell et al. (2026), extended to encompass the full MTZ and uppermost
lower mantle (Figure~\ref{fig:PYR-material-table}). We evaluate entropy
changes over a PT range of 1053--2673 K and 0.001--43 GPa using Gibbs
free energy minimization with Perple\_X (v7.0.9, Connolly, 2009),
assuming a dry pyrolitic bulk composition (Green et al., 1979) in the
Na\(_2\)O-CaO-FeO-MgO-Al\(_2\)O\(_3\)-SiO\(_2\) (NCFMAS) chemical system
with the equations of state and solution models of Stixrude \&
Lithgow-Bertelloni (2022). We determine an isentropic adiabat using the
Newton-Raphson method and evaluate reference material properties
(\(\bar{\rho}\), \(\bar{\alpha}\), \(\bar{C}_p\), \(\bar{\beta}\)) along
the adiabat using BurnMan (Cottaar et al., 2014; Myhill et al., 2023)
with equations of state from Stixrude \& Lithgow-Bertelloni (2022) for
pure Mg ringwoodite, bridgmanite, and periclase (Figure S1).

\begin{figure}
\centering
\includegraphics[width=0.7\linewidth,height=\textheight,keepaspectratio,alt={Entropy
(left) and density (right) changes across Earth's MTZ under
thermodynamic equilibrium and hydrostatic stress conditions. The black
box indicates the approximate PT range within our ASPECT simulations.
The white line indicates the isentropic adiabat used to calculate
reference material properties.}]{../figs/adiabat/PYR-material-table.png}
\caption{Entropy (left) and density (right) changes across Earth's MTZ
under thermodynamic equilibrium and hydrostatic stress conditions. The
black box indicates the approximate PT range within our ASPECT
simulations. The white line indicates the isentropic adiabat used to
calculate reference material properties.}\label{fig:PYR-material-table}
\end{figure}

\subsubsection{Initialization and Boundary
Conditions}\label{sec:initialization-and-boundary-conditions}

We conduct simulations in a 1500 \(\times\) 1000 km rectangular model
domain initialized with equilibrium mineral assemblages across the
olivine, wadsleyite, ringwoodite, and post-spinel stability fields. An
applied surface pressure of 10 GPa and surface temperature of 1706 K
places the 410, 520, and 660 transitions at approximately the center of
the model domain (Figure~\ref{fig:initial-setup-diagram}). We compute
initial adiabatic PT profiles by numerically integrating:

\begin{equation}\protect\phantomsection\label{eq:adiabatic-temperature}{\frac{d\bar{T}}{dy} = \frac{\bar{\alpha}\, \bar{T}\, g}{\bar{C}_p}}\end{equation}

\begin{equation}\protect\phantomsection\label{eq:adiabatic-pressure}{\frac{d\bar{P}}{dy} = \bar{\rho}\, g}\end{equation}

where the material properties \(\bar{\rho}\), \(\bar{\alpha}\), and
\(\bar{C}_p\) are determined from the stable equilibrium assemblage
evaluated along the adiabatic reference conditions (Figure S1).

We superimpose Gaussian thermal anomalies of -500 K (SW suite) or -1000
K (NC suite) on the adiabatic profile using linear features with 15 km
half-width Gaussian cross-sections and 5 km tanh-tapered ends. Slab
anomalies extend 105 km horizontally and 175 km downward from the top
boundary. We prescribe constant inflow velocities
\(\vec{u}_\mathrm{in}\) of 5 or 10 cm yr\(^{-1}\) parallel to the
thermal anomalies (plus 20 cm yr\(^{-1}\) in two sluggish-410
sensitivity tests;
Section~\ref{sec:numerical-experiment-design-and-postprocessing}).
Boundary conditions on lateral walls apply zero horizontal velocity with
a lithostatic normal traction, and the (open) top and bottom boundaries
apply constant normal traction equal to the initial lithostatic
pressure.

\begin{figure}
\centering
\includegraphics[width=0.5\linewidth,height=\textheight,keepaspectratio,alt={Initial
setup for slab simulations follows Kerswell et al. (2026), but with an
extended 1500 \textbackslash times 1000 km model domain and additional
phase transitions (410, 520, and 660; gray bands). Slab anomalies (-500
K in the SW suite shown here, or -1000 K in the NC suite) extend 105 km
horizontally and 175 km downward from the top boundary. All other
boundary conditions are consistent with Kerswell et al.
(2026).}]{../figs/setup/initial-setup-diagram.png}
\caption{Initial setup for slab simulations follows Kerswell et al.
(2026), but with an extended 1500 \(\times\) 1000 km model domain and
additional phase transitions (410, 520, and 660; gray bands). Slab
anomalies (-500 K in the SW suite shown here, or -1000 K in the NC
suite) extend 105 km horizontally and 175 km downward from the top
boundary. All other boundary conditions are consistent with Kerswell et
al. (2026).}\label{fig:initial-setup-diagram}
\end{figure}

\subsubsection{Material Model}\label{sec:material-model}

\paragraph{Material Properties}\label{sec:material-properties}

We compute material properties with BurnMan using equations of state
from Stixrude \& Lithgow-Bertelloni (2022) for pure Mg ringwoodite,
bridgmanite, and periclase (see
Section~\ref{sec:adiabatic-reference-conditions}), and all properties
except density are held at their adiabatic reference values (Figure S1).
We correct density using a first-order Taylor expansion:

\begin{equation}\protect\phantomsection\label{eq:density-ala}{\rho \approx \bar{\rho} \left(1 + \bar{\beta}\, \hat{P} - \bar{\alpha}\, \hat{T} \right)}\end{equation}

where \(\hat{P} = P - \bar{P}\) and \(\hat{T} = T - \bar{T}\) are the
dynamic pressure and temperature. The volume fraction fields
\(X_\mathrm{wd}\), \(X_\mathrm{ri}\), and \(X_\mathrm{ps}\) account for
phase-transition density jumps and evolve according to the kinetic
models described in Section~\ref{sec:reaction-kinetics}. We hold thermal
conductivity \(\bar{k}\) = 4.0 W m\(^{-1}\)K\(^{-1}\) constant in all
numerical experiments.

\paragraph{Reaction Kinetics}\label{sec:reaction-kinetics}

We treat reaction rates for all three MTZ phase transitions using the
site-saturated nucleation framework of Cahn (1956):

\begin{equation}\protect\phantomsection\label{eq:volume-fraction}{X = 1 - \exp\!\left(-N\, \dot{x}\, t \right)}\end{equation}

\begin{equation}\protect\phantomsection\label{eq:volume-transformation-rate}{\dot{X} = N\, \dot{x}\, \left(1 - X \right)}\end{equation}

where \(X\) is the volume fraction of the product phase, \(\dot{X}\) is
the volume transformation rate, \(N = \text{6.67}\, /\, d\) m\(^{-1}\)
is a geometric factor accounting for grain-boundary nucleation sites
with grain size \(d\), \(\dot{x}\) is the phase-specific growth rate,
and \(t\) is elapsed time after site saturation.

We formulate the 410 and 520 transitions using the interface-controlled
growth model of Hosoya et al. (2005), with macroscopic reaction rates
governed by:

\begin{equation}\protect\phantomsection\label{eq:reaction-rate-410}{\dot{X} = Z_\mathrm{ol}\, T\, \exp\!\left(-\frac{H^{\ast} + P V^{\ast}}{R\, T}\right) \left(1 - \exp\!\left[-\frac{\Delta G}{R\, T}\right] \right)\, \left(1 - X \right)}\end{equation}

where \(H^{\ast}\) = 274 kJ mol\(^{-1}\) and \(V^{\ast}\) = 3.0e-6
m\(^3\) mol\(^{-1}\). We hold the kinetic prefactors fixed at their
nominal value of \(Z_\mathrm{ol}\) = \(Z_\mathrm{wd}\) = 1.4e4 K
s\(^{-1}\) in the NC suite and at \(Z_\mathrm{ol}\) = \(Z_\mathrm{wd}\)
= 4.44e-4 K s\(^{-1}\) in the SW suite, which reflect intermediate 410
kinetics (Kerswell et al., 2026). At their respective slab temperatures
(-1000 K and -500 K inflow anomalies), both choices produce relatively
fast, near-equilibrium 410 and 520 reaction rates (Kerswell et al.,
2026) so that shallower transitions do not mask the kinetic signal
investigated at the 660.

The molar Gibbs free energy difference for each phase transformation is
approximated by:

\begin{equation}\protect\phantomsection\label{eq:gibbs}{\Delta G \approx \Delta \bar{G} + \hat{P}\, \Delta \bar{V} - \hat{T}\, \Delta \bar{S}}\end{equation}

where we compute the thermodynamic quantities \(\Delta \bar{G}\),
\(\Delta \bar{V}\), and \(\Delta \bar{S}\) for each major reaction along
the adiabatic reference profile (Figure S2).

The 660 post-spinel reaction requires a fundamentally different kinetic
formulation. Because one phase breaks down to two phases of different
compositions, diffusion must play a role regardless of whether interface
control has an effect. Microstructural observations show that the
forward ringwoodite \(\Rightarrow\) bridgmanite + periclase
transformation proceeds via diffusion-controlled eutectoid
decomposition, in which bridgmanite and periclase lamellae nucleate
rapidly at grain boundaries and grow at rates limited by silicon
self-diffusion (grain-boundary and volume) through the parent
ringwoodite (Dobson \& Mariani, 2014; Kubo et al., 2002; Yamazaki et
al., 2005; Zener, 1946). In classical and modern treatments of
diffusion-controlled lamellar growth (Burke, 1965; Christian, 2007;
Hillert, 1957, 1969; Porter \& Easterling, 2009; Zener, 1946), the
chemical-potential gradient driving diffusion across the advancing
reaction front scales with \(\Delta G / \lambda\), while the
interlamellar spacing \(\lambda\) that balances free-energy release
against interfacial energy scales inversely with overstepping
(\(\lambda \propto 1 / \Delta G\)). Combining these two factors yields a
growth rate proportional to \(\Delta G^2\):

\begin{equation}\protect\phantomsection\label{eq:growth-rate-660}{\dot{x} = K^{\prime}\, \exp\!\left(-\frac{E^{\ast}}{R\, T}\right)\, \frac{\Delta G^2}{\bar{V}_\mathrm{ps}^2}}\end{equation}

where \(K^{\prime}\) is the growth rate prefactor, \(E^{\ast}\) is the
activation energy for Si self-diffusion, and
\(\bar{V}_\mathrm{ps}^2 = (\bar{V}_\mathrm{bg} + \bar{V}_\mathrm{pe})^2\)
is the squared molar volume of the post-spinel assemblage. We refactor
\(\Delta G^2 = \Delta G \cdot |\Delta G|\) as the thermodynamic driving
force term, which preserves the sign of \(\Delta G\) for reverse
reactions when \(\Delta G\) \(<\) 0.

Combining \(N\), \(K^{\prime}\), and \(\bar{V}_\mathrm{ps}^2\) into a
single kinetic prefactor
\(Z_\mathrm{ri} = \text{6.67}\, K^{\prime}\, /\, d\, \bar{V}_\mathrm{ps}^2\),
the macroscopic reaction rate for ringwoodite \(\Leftrightarrow\)
bridgmanite + periclase becomes:

\begin{equation}\protect\phantomsection\label{eq:reaction-rate-660}{\dot{X} = Z_\mathrm{ri}\, \exp\!\left(-\frac{E^{\ast}}{R\, T}\right)\, \Delta G \cdot |\Delta G| \left(1 - X \right)}\end{equation}

Equation~\ref{eq:reaction-rate-660} introduces a fundamental shift from
the interface-controlled reaction mechanism at the 410 and 520
(Equation~\ref{eq:reaction-rate-410}) to a diffusion-controlled reaction
mechanism at the 660. Mathematically, the driving force at the 660 is
quadratic (\(\Delta G^2\)), whereas the 410 and 520 follow a linear
approximation near equilibrium (\(\Delta G\,/\,RT \approx\) 0) before
saturating at high overstepping. Physically, this quadratic dependence
more aggressively suppresses reaction rates near equilibrium for the 660
transition. Conversely, the lack of exponential saturation means the 660
reaction rate remains sensitive to the magnitude of overstepping even
far from equilibrium, unlike the 410 and 520 where the driving force
term trends towards unity. These divergent behaviors underpin the
distinct structural and dynamic signatures discussed in Sections
\ref{sec:results} and \ref{sec:discussion}.

We calibrate \(Z_\mathrm{ri}\) against the high-overpressure
experimental dataset of Kubo et al. (2002) using Monte Carlo regression
with a post-spinel molar volume of \(\bar{V}_\mathrm{ps}\) = 3.381e-5
m\(^3\) mol\(^{-1}\) (Figure S3). Our regression yields \(E^{\ast}\) =
310 \(\pm\) 233 kJ mol\(^{-1}\) and \(\ln(K^{\prime})\) = -29.8 \(\pm\)
22.1 (2\(\sigma\)), both consistent with Kubo et al. (2002) within
uncertainty. Propagating the \(\pm\) 2\(\sigma\) uncertainty in
\(\ln(K^{\prime})\) across typical mantle grain sizes of \(d \in\)
{[}1e-3, 10e-3{]} m (Karato, 1984, 2008) yields a 95\% confidence range
of \(Z_\mathrm{ri} \in\) {[}1.7e-11, 2.6e9{]} mol\(^2\) J\(^{-2}\)
s\(^{-1}\) (Figure S4). This wide 2\(\sigma\) envelope reflects both
systematic error (experimental uncertainties in controlling and
measuring temperature, pressure, and transformation progress under
extreme multi-anvil conditions) and random error (natural rock
variability arising from differences in grain size, defect structure,
and minor-element or volatile composition). Within this uncertainty
envelope, we fix \(E^{\ast}\) = 355 kJ mol\(^{-1}\) (after Kubo et al.,
2002) and sample two six-order-of-magnitude brackets (Figure S4). The NC
suite uses \(\ln(K^{\prime}) \in\) {[}-33.5, -22.0{]} and \(d \in\)
{[}1e-3, 10e-3{]} m to span \(Z_\mathrm{ri} \in\) {[}1.6e-3, 1.6e3{]}
mol\(^2\) J\(^{-2}\) s\(^{-1}\), while the SW suite spans
\(Z_\mathrm{ri} \in\) {[}5.08e-11, 5.08e-5{]} mol\(^2\) J\(^{-2}\)
s\(^{-1}\) (near the lower bound of the 95\% confidence interval in
Figure S4). Both suites therefore sample physically permitted kinetic
conditions ranging from slow decomposition in coarse-grained, dry rocks
to fast decomposition in fine-grained rocks.

\paragraph{Operator Splitting}\label{sec:operator-splitting}

Since reaction rates for the 410, 520, and 660 transitions can exceed
advective timescales, we employ the first-order operator splitting
scheme to decouple reaction from advection. Volume fractions are updated
in two sequential steps per advection timestep \(\Delta t\): 1) a
reaction substep integrating Equations \ref{eq:reaction-rate-410} and
\ref{eq:reaction-rate-660} over the full \(\Delta t\) using the ARKODE
adaptive Runge-Kutta integrator (Hindmarsh et al., 2005; Reynolds et
al., 2023) with relative tolerance 1e-6; and 2) an advection step:

\begin{equation}\protect\phantomsection\label{eq:composition}{\frac{\partial X}{\partial t} + \vec{u} \cdot \nabla X = 0}\end{equation}

The three kinetically active transitions require simultaneous
integration during the reaction step, with the ARKODE substep size
selected adaptively to satisfy the tolerance across all three reaction
systems.

\subsubsection{Rheological Model}\label{sec:rheological-model}

While the volume fractions \(X_\mathrm{ol}\), \(X_\mathrm{wd}\),
\(X_\mathrm{ri}\), and \(X_\mathrm{ps}\) track the actual mineralogy for
the purpose of density and buoyancy (including metastable phases), we
also define equilibrium fractions \(\phi_\mathrm{ol}\),
\(\phi_\mathrm{wd}\), \(\phi_\mathrm{ri}\), and \(\phi_\mathrm{ps}\).
These fractions represent the thermodynamic state along the reference
adiabat and facilitate smooth transitions in background viscosity. We
define a transition function \(f_j\) for each reaction \(j\):

\begin{equation}\protect\phantomsection\label{eq:transition-function}{f_j = \frac{1}{2} \left[ 1 - \tanh \left( \frac{\Delta G_j}{w} \right) \right]}\end{equation}

where \(w\) = 1000 J mol\(^{-1}\) is a ``Gibbs width'' that smooths the
equilibrium fractions \(\phi\) over a few kilometers, avoiding
discontinuous jumps in background viscosity at sharp equilibrium phase
boundaries. To ensure mass conservation \(\sum \phi_i = \text{1}\), we
partition the stability fractions as follows:

\begin{equation}\protect\phantomsection\label{eq:equilibrium-fractions}{\begin{aligned}
    \phi_\mathrm{ol} &= 1 - f_\mathrm{wd} \\
    \phi_\mathrm{wd} &= f_\mathrm{wd}\, (1 - f_\mathrm{ri}) \\
    \phi_\mathrm{ri} &= f_\mathrm{wd}\, f_\mathrm{ri}\, (1 - f_\mathrm{ps}) \\
    \phi_\mathrm{ps} &= f_\mathrm{wd}\, f_\mathrm{ri}\, f_\mathrm{ps}
  \end{aligned}}\end{equation}

The rheological model incorporates phase-dependent viscosity transitions
as a function of the equilibrium fractions \(\phi\) computed in
Equation~\ref{eq:equilibrium-fractions}. We use a temperature-dependent
Arrhenius viscosity modified by a phase-dependent viscosity jump
prefactor \(b\):

\begin{equation}\protect\phantomsection\label{eq:rheological-model}{\eta = b\, \bar{\eta}\, \exp\!\left(-B\,\frac{\hat{T}}{\bar{T}}\right)}\end{equation}

where \(\bar{\eta}\) = 1e21 Pa s is the background viscosity,
\(\bar{T}\) is the reference adiabatic temperature, and
\(\hat{T} = T - \bar{T}\) is the nonadiabatic temperature anomaly. We
vary the activation factor \(B\) between 5 and 10 to control thermal
sensitivity.

We calculate the viscosity jump prefactor \(b\) via a log-linear average
of per-phase viscosity prefactors
(\(\nu_\mathrm{ol}, \nu_\mathrm{wd}, \nu_\mathrm{ri}, \nu_\mathrm{ps}\)),
weighted by the equilibrium fractions \(\phi\) defined in
Equation~\ref{eq:equilibrium-fractions}:

\begin{equation}\protect\phantomsection\label{eq:rheological-jump-factor}{\ln(b) = \sum_{i} \phi_i \ln(\nu_i)}\end{equation}

We adopt this log-linear (geometric) average because it is commonly used
for viscosity contrasts across mineral phase mixtures and avoids the
sensitivity to the single weakest phase that a harmonic (Reuss) average
would introduce (Ji, 2004). By anchoring viscosity jumps to the
equilibrium fractions \(\phi\) rather than the actual (metastable)
mineralogy \(X\), the viscosity structure remains a function of the
adiabatic reference condition. This decoupling prevents spurious
viscosity patches arising from metastable phases within cold slabs,
allowing us to isolate the kinetic contribution to slab impedance via
buoyancy from rheological effects. To maintain numerical stability, we
impose a global viscosity floor and ceiling of 1e19 and 1e24 Pa s,
respectively.

\subsubsection{Numerical Stabilization of Dynamic Pressure
Oscillations}\label{sec:numerical-stabilization-of-dynamic-pressure-oscillations}

Following Gassmöller et al. (2020) and Kerswell et al. (2026), we
exclude the dynamic pressure \(\hat{P}\) contribution to \(\Delta G\) in
the kinetic rate equations (Equations \ref{eq:reaction-rate-410} and
\ref{eq:reaction-rate-660}) while retaining it in the density
formulation (Equation~\ref{eq:density-ala}). In practice, this means
applying \(\Delta G \approx \Delta \bar{G} - \hat{T}\, \Delta \bar{S}\)
to the kinetic rate equations, which preserves the primary thermal
control of the reactions while eliminating pressure-wave feedback. The
quadratic \(\Delta G \cdot |\Delta G|\) driving force amplifies
pressure-wave sensitivity relative to the exponential form used at the
410 and 520, making this stabilization particularly important for the
660. Limitations of this approximation are discussed in
Section~\ref{sec:uncertainties-and-model-limitations}.

\subsubsection{Numerical Experiment Design and
Postprocessing}\label{sec:numerical-experiment-design-and-postprocessing}

We explore the ringwoodite decomposition kinetic parameter space across
16 baseline scenario configurations spanning slab thermal anomalies
(\(\Delta T_0\) = -500 and -1000 K), inflow velocities
(\(\vec{u}_\mathrm{in}\) = 5 and 10 cm yr\(^{-1}\)), rheological
strength contrasts (\(B\) = 5 and 10), and lower-mantle viscosity jumps
(\(b\) = 1 and 50; Table~\ref{tbl:scenarios}). Within each of the eight
baseline kinematic-rheological scenarios per suite, we sample eight
values of \(Z_\mathrm{ri}\) across six orders of magnitude, running the
SW suite as scenarios 0--7 (64 simulations) and repeating the same
kinematic-rheological grid for the NC suite as scenarios 8--15 (64
simulations). We also include four additional test simulations
(scenarios 16--19, T1--T4; Table~\ref{tbl:scenarios} and Table S1) at
\(\Delta T_0\) = -1000 K, \(B\) = 5, \(b\) = 1, \(\vec{u}_\mathrm{in}\)
= 5 and 20 cm yr\(^{-1}\), \(Z_\mathrm{ri}\) = 6.0e-1 mol\(^2\)
J\(^{-2}\) s\(^{-1}\), and \(Z_\mathrm{ol}\) = \(Z_\mathrm{wd}\) = 3.0
and 7.0 K s\(^{-1}\). As expected, these test simulations reproduced
analogous metastable olivine wedges shown in previous studies of very
cold slabs (Däßler et al., 1996; Devaux et al., 1997; Hosoya et al.,
2005; Riedel \& Karato, 1997; Rubie \& Ross, 1994).

We characterize 660 structure using two diagnostic measures extracted
along a representative depth profile that tracks the slab's
thermal-anomaly core across the transition zone in each simulation.
`Displacement' (\(\delta z\)) is the difference between the depth at
which the post-spinel volume fraction \(X_\mathrm{ps}\) reaches 0.9 and
the nominal equilibrium post-spinel depth (614.7 km along the reference
adiabat). `Width' (\(\delta w\)) is the depth interval between
\(X_\mathrm{ps}\) = 0.9 and \(X_\mathrm{ps}\) = 0.1. We also report the
maximum post-spinel reaction rate \(\dot{X}_\mathrm{max}\) and the
maximum descent velocity \(\vec{v}_\mathrm{max}\) along the same
profile, extracted within the phase transition zone (defined by the
measured 660 width \(\delta w\)). From these two quantities, we define a
characteristic reaction timescale (in Ma):

\begin{equation}\protect\phantomsection\label{eq:reaction-timescale}{\tau = \frac{1}{\dot{X}_\mathrm{max}}}\end{equation}

and a characteristic reaction completion length scale (in km):

\begin{equation}\protect\phantomsection\label{eq:characteristic-length-scale}{L = 10\, \frac{\vec{v}_\mathrm{max}}{\dot{X}_\mathrm{max}}}\end{equation}

which approximates the depth traversed during one characteristic
reaction timescale \(\tau\). \(L\) and the measured 660 width
\(\delta w\) are related but not identical. \(L\) is a local estimate
based on the maximum reaction rate and maximum vertical velocity through
the 660, while \(\delta w\) measures the full spatial extent of the
transition between \(X_\mathrm{ps}\) = 0.1 and 0.9. The two metrics
diverge where reaction rate varies strongly within the phase transition
zone.

Because the two suites differ in both imposed thermal anomaly and
kinetic prefactor, total displacement \(\delta z\) combines thermal
displacement with kinetic inhibition. To separate these two
contributions at both the 660 and the 410 (nominal reference depth 364.2
km), we interpolate the dynamic temperature \(\hat{T}\) along the
representative depth profile at the measured \(X\) = 0.9 transition
depth and compute the equilibrium thermal displacement contribution:

\begin{equation}\protect\phantomsection\label{eq:thermal-displacement}{\delta z_t = \frac{\gamma\, \hat{T}}{\rho\, g}}\end{equation}

and the excess kinetic displacement contribution:

\begin{equation}\protect\phantomsection\label{eq:kinetic-displacement}{\delta z_k = \delta z - \delta z_t}\end{equation}

using \(\rho\) = 4100 kg m\(^{-3}\), \(g\) = 9.81 m s\(^{-2}\), and the
pure-Mg reference Clapeyron slopes \(\gamma_{660}\) = -1.869 MPa
K\(^{-1}\) and \(\gamma_{410}\) = +3.090 MPa K\(^{-1}\).

We use \(L\) and \(\vec{v}_\mathrm{max}\) to heuristically classify each
slab simulation into three regimes using fixed threshold values. The
quasi-equilibrium regime corresponds to \(L\) \(\leq\) 5 km and
\(\vec{v}_\mathrm{max}\) \(>\) 0.5 cm yr\(^{-1}\); the intermediate
regime corresponds to \(L\) \(>\) 5 km and \(\vec{v}_\mathrm{max}\)
\(>\) 0.5 cm yr\(^{-1}\); and the stagnation regime corresponds to
\(\vec{v}_\mathrm{max}\) \(\leq\) 0.5 cm yr\(^{-1}\).

\begin{longtable}[]{@{}
  >{\raggedright\arraybackslash}p{(\linewidth - 16\tabcolsep) * \real{0.0417}}
  >{\raggedleft\arraybackslash}p{(\linewidth - 16\tabcolsep) * \real{0.0667}}
  >{\raggedleft\arraybackslash}p{(\linewidth - 16\tabcolsep) * \real{0.1000}}
  >{\raggedleft\arraybackslash}p{(\linewidth - 16\tabcolsep) * \real{0.1750}}
  >{\raggedleft\arraybackslash}p{(\linewidth - 16\tabcolsep) * \real{0.0250}}
  >{\raggedleft\arraybackslash}p{(\linewidth - 16\tabcolsep) * \real{0.0250}}
  >{\raggedleft\arraybackslash}p{(\linewidth - 16\tabcolsep) * \real{0.2417}}
  >{\raggedleft\arraybackslash}p{(\linewidth - 16\tabcolsep) * \real{0.1583}}
  >{\raggedleft\arraybackslash}p{(\linewidth - 16\tabcolsep) * \real{0.1667}}@{}}
\caption{Thermal, kinematic, rheological, and kinetic parameters across
all simulations. T1--T4 are sluggish-410 sensitivity tests in cold slabs
(not shown). Units are \(\Delta T_0\): K; \(\vec{u}_\mathrm{in}\): cm
yr\(^{-1}\); \(B\): dimensionless; \(b\): dimensionless;
\(Z_\mathrm{ol}, Z_\mathrm{wd}\): K s\(^{-1}\); \(Z_\mathrm{ri}\):
mol\(^2\) J\(^{-2}\) s\(^{-1}\).}\label{tbl:scenarios}\tabularnewline
\toprule\noalign{}
\begin{minipage}[b]{\linewidth}\raggedright
Suite
\end{minipage} & \begin{minipage}[b]{\linewidth}\raggedleft
Scenario
\end{minipage} & \begin{minipage}[b]{\linewidth}\raggedleft
\(\Delta T_0\)
\end{minipage} & \begin{minipage}[b]{\linewidth}\raggedleft
\(\vec{u}_\mathrm{in}\)
\end{minipage} & \begin{minipage}[b]{\linewidth}\raggedleft
\(B\)
\end{minipage} & \begin{minipage}[b]{\linewidth}\raggedleft
\(b\)
\end{minipage} & \begin{minipage}[b]{\linewidth}\raggedleft
\(Z_\mathrm{ol},Z_\mathrm{wd}\)
\end{minipage} & \begin{minipage}[b]{\linewidth}\raggedleft
low \(Z_\mathrm{ri}\)
\end{minipage} & \begin{minipage}[b]{\linewidth}\raggedleft
high \(Z_\mathrm{ri}\)
\end{minipage} \\
\midrule\noalign{}
\endfirsthead
\toprule\noalign{}
\begin{minipage}[b]{\linewidth}\raggedright
Suite
\end{minipage} & \begin{minipage}[b]{\linewidth}\raggedleft
Scenario
\end{minipage} & \begin{minipage}[b]{\linewidth}\raggedleft
\(\Delta T_0\)
\end{minipage} & \begin{minipage}[b]{\linewidth}\raggedleft
\(\vec{u}_\mathrm{in}\)
\end{minipage} & \begin{minipage}[b]{\linewidth}\raggedleft
\(B\)
\end{minipage} & \begin{minipage}[b]{\linewidth}\raggedleft
\(b\)
\end{minipage} & \begin{minipage}[b]{\linewidth}\raggedleft
\(Z_\mathrm{ol},Z_\mathrm{wd}\)
\end{minipage} & \begin{minipage}[b]{\linewidth}\raggedleft
low \(Z_\mathrm{ri}\)
\end{minipage} & \begin{minipage}[b]{\linewidth}\raggedleft
high \(Z_\mathrm{ri}\)
\end{minipage} \\
\midrule\noalign{}
\endhead
\bottomrule\noalign{}
\endlastfoot
SW & 0 & -500 & 5 & 5 & 50 & 4.44e-4 & 5.08e-11 & 5.08e-5 \\
SW & 1 & -500 & 5 & 10 & 50 & 4.44e-4 & 5.08e-11 & 5.08e-5 \\
SW & 2 & -500 & 10 & 5 & 50 & 4.44e-4 & 5.08e-11 & 5.08e-5 \\
SW & 3 & -500 & 10 & 10 & 50 & 4.44e-4 & 5.08e-11 & 5.08e-5 \\
SW & 4 & -500 & 5 & 5 & 1 & 4.44e-4 & 5.08e-11 & 5.08e-5 \\
SW & 5 & -500 & 5 & 10 & 1 & 4.44e-4 & 5.08e-11 & 5.08e-5 \\
SW & 6 & -500 & 10 & 5 & 1 & 4.44e-4 & 5.08e-11 & 5.08e-5 \\
SW & 7 & -500 & 10 & 10 & 1 & 4.44e-4 & 5.08e-11 & 5.08e-5 \\
NC & 8 & -1000 & 5 & 5 & 50 & 1.4e4 & 1.6e-3 & 1.6e3 \\
NC & 9 & -1000 & 5 & 10 & 50 & 1.4e4 & 1.6e-3 & 1.6e3 \\
NC & 10 & -1000 & 10 & 5 & 50 & 1.4e4 & 1.6e-3 & 1.6e3 \\
NC & 11 & -1000 & 10 & 10 & 50 & 1.4e4 & 1.6e-3 & 1.6e3 \\
NC & 12 & -1000 & 5 & 5 & 1 & 1.4e4 & 1.6e-3 & 1.6e3 \\
NC & 13 & -1000 & 5 & 10 & 1 & 1.4e4 & 1.6e-3 & 1.6e3 \\
NC & 14 & -1000 & 10 & 5 & 1 & 1.4e4 & 1.6e-3 & 1.6e3 \\
NC & 15 & -1000 & 10 & 10 & 1 & 1.4e4 & 1.6e-3 & 1.6e3 \\
T1 & 16 & -1000 & 5 & 5 & 1 & 3.0e0 & 6.0e-1 & \\
T2 & 17 & -1000 & 5 & 5 & 1 & 7.0e0 & 6.0e-1 & \\
T3 & 18 & -1000 & 20 & 5 & 1 & 3.0e0 & 6.0e-1 & \\
T4 & 19 & -1000 & 20 & 5 & 1 & 7.0e0 & 6.0e-1 & \\
\end{longtable}

\section{Results}\label{sec:results}

\subsection{Simulation Snapshots}\label{sec:simulation-snapshots}

Representative slab snapshots after 100 Ma demonstrate that sluggish
ringwoodite kinetics (SW suite) drive slabs monotonically from
penetration to stagnation as \(Z_\mathrm{ri}\) decreases, whereas
nominal ringwoodite kinetics (NC suite) hold the 660 near its
equilibrium depth across all six orders of magnitude in
\(Z_\mathrm{ri}\) (Figures \ref{fig:slab-composition-scen04} and
\ref{fig:slab-composition-scen12}; scenarios 4 and 12 in
Table~\ref{tbl:scenarios}; see Supplementary Information for additional
representative scenarios).

In the SW suite (Figure~\ref{fig:slab-composition-scen04}), three
kinetic regimes emerge with decreasing \(Z_\mathrm{ri}\). In the
quasi-equilibrium regime (\(Z_\mathrm{ri}\) = 5.08e-5 mol\(^2\)
J\(^{-2}\) s\(^{-1}\)), the slab descends through the 660 at 1.55 cm
yr\(^{-1}\) and the 660 remains narrow (\(\delta w\) = 6.3 km) and
modestly displaced (\(\delta z\) = 25.1 km, comprising \(\delta z_t\) =
16.5 km and \(\delta z_k\) = 8.6 km). In the intermediate regime
(\(Z_\mathrm{ri}\) = 1.90e-8 mol\(^2\) J\(^{-2}\) s\(^{-1}\)),
metastable ringwoodite accumulates and slows descent through the 660 to
1.11 cm yr\(^{-1}\), deepening the 660 by \(\delta z\) = 39.6 km
(\(\delta z_t\) = 17.1 km, \(\delta z_k\) = 22.4 km) and widening it to
\(\delta w\) = 11.8 km. In the stagnation regime (\(Z_\mathrm{ri}\) =
5.08e-11 mol\(^2\) J\(^{-2}\) s\(^{-1}\)), persistent metastable
buoyancy arrests descent (\(\vec{v}_\mathrm{max}\) \(\leq\) 0.25 cm
yr\(^{-1}\)) and the slab ponds above the 660, displacing it by
\(\delta z\) = 48.4 km (\(\delta z_t\) = 11.0 km, \(\delta z_k\) = 37.4
km) and broadening it to \(\delta w\) = 10.8 km.

In the NC suite (Figure~\ref{fig:slab-composition-scen12}), the slab
penetrates continuously into the lower mantle and the 660 remains narrow
across the entire six-order-of-magnitude range of \(Z_\mathrm{ri}\).
Even at the slowest nominal prefactor (\(Z_\mathrm{ri}\) = 1.6e-3
mol\(^2\) J\(^{-2}\) s\(^{-1}\)), the slab descends through the 660 at
2.96 cm yr\(^{-1}\) with a \(\delta w\) = 6.3 km wide 660 displaced
\(\delta z\) = 40.1 km downward, of which \(\delta z_t\) = 29.6 km is
thermal displacement and only \(\delta z_k\) = 10.5 km is kinetic
inhibition. Under intermediate and fast nominal kinetics
(\(Z_\mathrm{ri}\) = 6.0e-1 and 1.6e3 mol\(^2\) J\(^{-2}\) s\(^{-1}\)),
the slab descends at 3.04--3.05 cm yr\(^{-1}\) with a \(\delta w\) = 8.1
km wide 660 displaced by \(\delta z\) = 37.4 km (\(\delta z_t\) = 29.6
km, \(\delta z_k\) = 7.8 km). At nominal \(Z_\mathrm{ri}\) values from
Kubo et al. (2002), large overstepping in the slab drives rapid
quadratic reaction rates (Equation~\ref{eq:reaction-rate-660}) that
prevent an extended metastable ringwoodite zone from accumulating, even
in very cold slab cores (\(\hat{T} \approx\) -358 to -839 K at 660 km).
This contrasts with the stagnant ponding observed in the SW suite, which
is produced by very sluggish kinetics near the lower limit of the
2\(\sigma\) experimental uncertainties (Figures S3--S4).

\clearpage

\begin{figure}
\centering
\includegraphics[width=1\linewidth,height=\textheight,keepaspectratio,alt={Reference
SW suite snapshots after 100 Ma with moderate inflow velocity,
intermediate slab strength, and no lower-mantle viscosity jump (scenario
4; Table~). Simulations show stagnation (top row), intermediate (middle
row), and quasi-equilibrium (bottom row) kinetic regimes. Columns show
dynamic temperature \textbackslash hat\{T\} (left column), dynamic
density \textbackslash hat\{\textbackslash rho\} (middle column), and
vertical velocity \textbackslash vec\{u\}\_y (right column). Thick black
lines (left column) highlight the 10\% and 90\% wadsleyite and
post-spinel volume fraction contours used to define 410 and 660
displacement (\textbackslash delta z) and width (\textbackslash delta
w). The background grid indicates ASPECT's adaptive mesh refinement.
White trace (left column) indicates the representative depth profile
used to extract data, and the white arrows indicate where transition
structure is
evaluated.}]{../figs/simulation/compositions/slab-set1-composition-scen04-0100.png}
\caption{Reference SW suite snapshots after 100 Ma with moderate inflow
velocity, intermediate slab strength, and no lower-mantle viscosity jump
(scenario 4; Table~\ref{tbl:scenarios}). Simulations show stagnation
(top row), intermediate (middle row), and quasi-equilibrium (bottom row)
kinetic regimes. Columns show dynamic temperature \(\hat{T}\) (left
column), dynamic density \(\hat{\rho}\) (middle column), and vertical
velocity \(\vec{u}_y\) (right column). Thick black lines (left column)
highlight the 10\% and 90\% wadsleyite and post-spinel volume fraction
contours used to define 410 and 660 displacement (\(\delta z\)) and
width (\(\delta w\)). The background grid indicates ASPECT's adaptive
mesh refinement. White trace (left column) indicates the representative
depth profile used to extract data, and the white arrows indicate where
transition structure is evaluated.}\label{fig:slab-composition-scen04}
\end{figure}

\clearpage

\begin{figure}
\centering
\includegraphics[width=1\linewidth,height=\textheight,keepaspectratio,alt={Reference
NC suite snapshots after 100 Ma at identical kinematic and rheological
settings (scenario 12; Table~). Simulations span slow (top row),
intermediate (middle row), and fast (bottom row) nominal kinetic
prefactors. All other annotations follow Figure~. Slabs penetrate
continuously into the lower mantle with narrow (\textbackslash delta w =
6.3--8.1 km), thermally deepened (\textbackslash delta z\_t = 29.6 km,
\textbackslash delta z\_k = 7.8--10.5 km) 660s across the entire nominal
Z\_\textbackslash mathrm\{ri\}
range.}]{../figs/simulation/compositions/slab-set1-composition-scen12-0100.png}
\caption{Reference NC suite snapshots after 100 Ma at identical
kinematic and rheological settings (scenario 12;
Table~\ref{tbl:scenarios}). Simulations span slow (top row),
intermediate (middle row), and fast (bottom row) nominal kinetic
prefactors. All other annotations follow
Figure~\ref{fig:slab-composition-scen04}. Slabs penetrate continuously
into the lower mantle with narrow (\(\delta w\) = 6.3--8.1 km),
thermally deepened (\(\delta z_t\) = 29.6 km, \(\delta z_k\) = 7.8--10.5
km) 660s across the entire nominal \(Z_\mathrm{ri}\)
range.}\label{fig:slab-composition-scen12}
\end{figure}

\clearpage

\subsection{Structure of the 660: Displacement, Width, and
Thermal-Kinetic Contributions}\label{sec:structure-of-the-660}

Across all 128 slab simulations, total 660 displacements overlap between
the two suites (median \(\delta z\) = 38.0 km for SW versus 34.2 km for
NC), but decomposing \(\delta z\) into its kinetic (\(\delta z_k\);
middle panel) and thermal (\(\delta z_t\); right panel) contributions
reveals that the SW suite is dominated by progressive kinetic inhibition
while the NC suite is dominated by thermal displacement
(Figure~\ref{fig:contributions}).

Evaluating the characteristic reaction completion length scale \(L\)
(Equation~\ref{eq:characteristic-length-scale}) against descent velocity
through the 660 (\(\vec{v}_\mathrm{max}\)) separates the simulations
into three kinetic regimes bounded by \(L\) = 5 km and
\(\vec{v}_\mathrm{max}\) = 0.5 cm yr\(^{-1}\)
(Figure~\ref{fig:contributions}, left panel; see
Section~\ref{sec:numerical-experiment-design-and-postprocessing}). The
SW suite spans all three regimes (20 simulations in quasi-equilibrium,
18 in intermediate, and 26 in stagnation; Table S1). In the
quasi-equilibrium regime (region 1 in Figure~\ref{fig:contributions},
left panel), the characteristic reaction timescale \(\tau\)
(Equation~\ref{eq:reaction-timescale}) remains below 2 Ma and slabs
descend through the 660 at 0.5--3.2 cm yr\(^{-1}\), with the 660
displaced \(\delta z\) = 22--33 km downward (\(\delta z_k\) = 6.7--13.6
km) and \(\delta w\) = 4--12 km wide. In the intermediate regime (region
2 in Figure~\ref{fig:contributions}, left panel), \(\tau\) spans
0.3--5.1 Ma and descent velocities through the 660 of 0.6--2.9 cm
yr\(^{-1}\) produce continuous but progressively slower slab
penetration, with the 660 deepening by \(\delta z\) = 30--80 km
(\(\delta z_k\) = 12.4--61.4 km) and widening to \(\delta w\) = 7--34
km. In the stagnation regime (region 3 in
Figure~\ref{fig:contributions}, left panel), descent velocities through
the 660 of 0.1--0.5 cm yr\(^{-1}\) and \(\tau\) up to 26 Ma leave the
660 displaced \(\delta z\) = 27--80 km downward (\(\delta z_k\) =
8.7--64.1 km) with widths of \(\delta w\) = 6--34 km. By contrast, 62 of
the 64 NC suite simulations remain in the quasi-equilibrium regime
(region 1) with \(L \ll\) 5 km and \(\vec{v}_\mathrm{max}\) = 0.9--6.6
cm yr\(^{-1}\). Only two NC simulations (scenario 14 at
\(Z_\mathrm{ri}\) = 1.6e-3--1.2e-2 mol\(^2\) J\(^{-2}\) s\(^{-1}\))
enter the intermediate regime, and none stagnate.

Decomposing 660 displacement into \(\delta z_k\) and \(\delta z_t\)
(Figure~\ref{fig:contributions}, middle and right panels) shows that
sluggish kinetics amplify thermal displacements by two to five times
without an abrupt re-sharpening threshold. In the SW suite (red points),
warmer slab-core thermal anomalies at the 660 (\(\hat{T} \approx\) -236
to -457 K) produce a narrow band of thermal contributions spanning
\(\delta z_t\) = 11.0--21.2 km (median 16.9 km). As
\(\dot{X}_\mathrm{max}\) decreases from 1.35e2 to 3.81e-2 Ma\(^{-1}\),
the kinetic contribution progresses from \(\delta z_k\) = 6.7--10.7 km
at the fastest \(Z_\mathrm{ri}\) to 37.4--64.1 km at the slowest
\(Z_\mathrm{ri}\) (suite median 18.2 km), while 660 width progresses
from \(\delta w\) = 4.0--12.0 to 10.8--33.7 km. In contrast to the
abrupt re-sharpening threshold at the 410 (Kerswell et al., 2026),
\(\delta z_k\) and \(\delta w\) at the 660 change monotonically with
decreasing reaction rate across all three regimes, with no re-sharpening
at any tested kinetic prefactor. In the NC suite (blue points), colder
slab cores at the 660 (\(\hat{T} \approx\) -358 to -839 K) generate
larger thermal contributions of \(\delta z_t\) = 16.6--39.0 km (median
25.5 km), while \(\delta z_k\) remains flat at 5.6--12.5 km (median 8.2
km) with narrow 660 widths of \(\delta w\) = 3.2--13.8 km (median 6.5
km) across nearly the entire six-order-of-magnitude range of
\(Z_\mathrm{ri}\). Both suites collapse onto a single continuous
\(\delta z_k\) curve (Figure~\ref{fig:contributions}, middle panel)
where the kinetically fastest SW simulations meet the NC
quasi-equilibrium baseline near \(\delta z_k \approx\) 6--11 km. Only
the kinematically fastest, rheologically weakest NC configuration at the
slowest nominal \(Z_\mathrm{ri}\) (scenario 14; \(\delta z_k\) =
15.5--22.3 km) departs from this baseline to overlap the SW suite trend.
These two cases are the only NC simulations where rapid descent outpaces
ringwoodite transformation.

\begin{figure}
\centering
\includegraphics[width=1\linewidth,height=\textheight,keepaspectratio,alt={Kinetic
regime classification (left), kinetic contribution \textbackslash delta
z\_k (middle), and thermal contribution \textbackslash delta z\_t
(right) to 660 displacement across all 128 slab simulations after 100 Ma
(Table S1). Point size scales with 660 width \textbackslash delta w
(3.2--34.3 km), and color distinguishes the SW suite (red) and the NC
suite (blue). Left: maximum descent velocity through the 660
(\textbackslash vec\{v\}\_\textbackslash mathrm\{max\}) versus
characteristic reaction length scale L, with dashed lines at L = 5 km
and \textbackslash vec\{v\}\_\textbackslash mathrm\{max\} = 0.5 cm
yr\^{}\{-1\} dividing (1) quasi-equilibrium, (2) intermediate, and (3)
stagnation regimes. Middle and right: \textbackslash delta z\_k and
\textbackslash delta z\_t versus maximum reaction rate \textbackslash
dot\{X\}\_\textbackslash mathrm\{max\}, showing that both suites
collapse onto a single \textbackslash delta z\_k curve while
\textbackslash delta z\_t separates into distinct bands governed by
slab-core temperature and \textbackslash vec\{v\}\_\textbackslash
mathrm\{max\} at the 660.}]{../figs/structure/contributions.png}
\caption{Kinetic regime classification (left), kinetic contribution
\(\delta z_k\) (middle), and thermal contribution \(\delta z_t\) (right)
to 660 displacement across all 128 slab simulations after 100 Ma (Table
S1). Point size scales with 660 width \(\delta w\) (3.2--34.3 km), and
color distinguishes the SW suite (red) and the NC suite (blue). Left:
maximum descent velocity through the 660 (\(\vec{v}_\mathrm{max}\))
versus characteristic reaction length scale \(L\), with dashed lines at
\(L\) = 5 km and \(\vec{v}_\mathrm{max}\) = 0.5 cm yr\(^{-1}\) dividing
(1) quasi-equilibrium, (2) intermediate, and (3) stagnation regimes.
Middle and right: \(\delta z_k\) and \(\delta z_t\) versus maximum
reaction rate \(\dot{X}_\mathrm{max}\), showing that both suites
collapse onto a single \(\delta z_k\) curve while \(\delta z_t\)
separates into distinct bands governed by slab-core temperature and
\(\vec{v}_\mathrm{max}\) at the 660.}\label{fig:contributions}
\end{figure}

\clearpage

\subsection{Interactions Among the Key Physical
Parameters}\label{sec:interactions-among-key-physical-parameters}

Reaction kinetics interact with lower-mantle viscosity jump \(b\),
inflow velocity \(\vec{u}_\mathrm{in}\), and slab strength \(B\) to
govern where slabs stagnate and how deeply metastable ringwoodite is
advected (Figure~\ref{fig:interactions}; Table S1).

Within the SW suite (Figure~\ref{fig:interactions}, left column),
decreasing \(Z_\mathrm{ri}\) from 5.08e-5 to 5.08e-11 mol\(^2\)
J\(^{-2}\) s\(^{-1}\) reduces \(\dot{X}_\mathrm{max}\) proportionally
from 1.35e2 to 3.81e-2 Ma\(^{-1}\) (Table S1;
Equation~\ref{eq:reaction-rate-660}), progressively slowing descent
through the 660 (top left), amplifying \(\delta z_k\) (second row, left)
and \(\delta w\) (bottom left), and slightly reducing \(\delta z_t\) at
the lowest \(Z_\mathrm{ri}\) as stagnating slabs warm by conduction
(third row, left). Within the NC suite (Figure~\ref{fig:interactions},
right column), decreasing \(Z_\mathrm{ri}\) from 1.6e3 to 1.6e-3
mol\(^2\) J\(^{-2}\) s\(^{-1}\) similarly reduces
\(\dot{X}_\mathrm{max}\) from 2.27e9 to 3.72e0 Ma\(^{-1}\), but because
reaction timescales remain short (\(\tau \ll\) 1 Ma), descent rate,
\(\delta z_k\), \(\delta z_t\), and \(\delta w\) remain nearly invariant
with \(Z_\mathrm{ri}\) and stratify instead by kinematic and rheological
settings.

Comparing dashed (\(b\) = 50) and solid (\(b\) = 1) curves of matching
color in Figure~\ref{fig:interactions} (see also Figures S6--S7 and
S13--S14) shows that a 50-fold lower-mantle viscosity jump slows slab
descent through the 660 in both suites (top row), but triggers
stagnation only when paired with sluggish kinetics (SW suite; top left).
In the NC suite (top right), imposing the same viscosity jump reduces
\(\vec{v}_\mathrm{max}\) from 2.9--6.6 cm yr\(^{-1}\) to 0.9--1.9 cm
yr\(^{-1}\) and folds the slab above the lower mantle (Figures S13--S14
and S17--S18), yet all eight NC scenarios remain above the 0.5 cm
yr\(^{-1}\) stagnation threshold (gray band), with flat \(\delta z_k\)
across the \(b\) = 50 curves (5.6--12.5 km; second row, right). In the
SW suite, where metastable ringwoodite buoyancy already resists descent,
the same viscosity jump drives the curves into the gray stagnation zone
at up to four orders of magnitude higher \(Z_\mathrm{ri}\) (top left).
This shift is largest for weak, slow slabs (light-blue circles).
Scenario 4 (\(b\) = 1, solid light-blue curve) crosses into stagnation
at or below \(Z_\mathrm{ri}\) = 2.63e-9 mol\(^2\) J\(^{-2}\) s\(^{-1}\),
whereas scenario 0 (\(b\) = 50, dashed light-blue curve; Figures S6--S7)
remains below 0.5 cm yr\(^{-1}\) across the entire tested
\(Z_\mathrm{ri}\) range. For stronger SW slabs (\(B\) = 10, dark-blue
squares and red triangles), the viscosity jump still reduces
\(\vec{v}_\mathrm{max}\) by a factor of two to three, but shifts the
stagnation crossing by only one order of magnitude or less in
\(Z_\mathrm{ri}\) (compare solid versus dashed curves).

Comparing fast-inflow (orange diamonds and red triangles) and
slow-inflow (light-blue circles and dark-blue squares) curves in
Figure~\ref{fig:interactions} demonstrates how advection modulates both
thermal and kinetic contributions to 660 structure (see Figures S8--S9
and S15--S16). In both suites, doubling \(\vec{u}_\mathrm{in}\) nearly
doubles \(\vec{v}_\mathrm{max}\) at \(b\) = 1 (top row, solid orange/red
versus solid blue curves) and transports colder slab material to
transition-zone depths. In the NC suite, faster inflow preserves colder
slab cores at the 660 (\(\hat{T}\) down to -839 K in scenario 14 versus
-637 K in scenario 12), increasing \(\delta z_t\) from 29.6 km to
37.8--39.0 km for weak slabs (third row, right, solid orange diamonds
versus solid light-blue circles). At the slowest nominal
\(Z_\mathrm{ri}\) (1.6e-3--1.2e-2 mol\(^2\) J\(^{-2}\) s\(^{-1}\)),
rapid descent in scenario 14 (\(\vec{v}_\mathrm{max}\) = 4.8--5.5 cm
yr\(^{-1}\); Figures S15--S16) begins to outpace ringwoodite
decomposition, producing the only upward inflection in NC \(\delta z_k\)
(to 15.5--22.3 km; second row, right) and \(\delta w\) (to 10.2--13.1
km; bottom right). In the SW suite, faster inflow shifts the stagnation
crossing to 1--2 orders of magnitude lower \(Z_\mathrm{ri}\) (top left)
while advecting untransformed ringwoodite deeper before reaction
completes, elevating \(\delta z_k\) to 7.9--64.1 km and \(\delta w\) to
4.0--34.3 km for \(\vec{u}_\mathrm{in}\) = 10 cm yr\(^{-1}\) (orange and
red curves; Figures S8--S11), compared to \(\delta z_k\) = 6.7--45.7 km
and \(\delta w\) = 4.0--21.4 km for \(\vec{u}_\mathrm{in}\) = 5 cm
yr\(^{-1}\) (blue curves; Figures \ref{fig:slab-composition-scen04},
S5--S7).

Finally, slab strength (dark-blue squares and red triangles versus
light-blue circles and orange diamonds) controls how coherently slabs
transmit buoyant and viscous stresses across the 660
(Figure~\ref{fig:interactions}; see Figures S10--S11 and S17--S18). In
the NC suite, stronger temperature-dependent viscosity (\(B\) = 10)
stiffens the outer thermal boundary layer of the slab, moderating core
descent velocity at \(b\) = 1 (4.5--4.6 cm yr\(^{-1}\) in scenario 15
versus 4.8--6.6 cm yr\(^{-1}\) in scenario 14; top right) and sampling
warmer temperatures along the 660 profile (\(\delta z_t\) = 16.6--21.8
km) than in weaker \(B\) = 5 slabs (\(\delta z_t\) = 29.1--39.0 km;
third row, right). In the SW suite, weaker slabs (\(B\) = 5) decelerate
more readily under metastable buoyancy and stagnate at up to two orders
of magnitude higher \(Z_\mathrm{ri}\) than stronger slabs (\(B\) = 10;
top left). Because early stagnation halts further downward advection of
ringwoodite, weaker SW slabs reach smaller maximum kinetic displacements
(\(\delta z_k\) = 37.4--58.5 km; total \(\delta z\) = 48.4--74.9 km) and
widths (\(\delta w\) = 10.8--24.7 km) at the lowest \(Z_\mathrm{ri}\).
Conversely, the strongest and fastest SW slabs (red triangles; Figures
S10--S11) resist stagnation down to \(Z_\mathrm{ri} \leq\) 3.81e-10
mol\(^2\) J\(^{-2}\) s\(^{-1}\), producing the largest kinetic
displacements (\(\delta z_k\) = 61.4--64.1 km; total \(\delta z\) =
78.3--80.5 km; second row, left) and widest transition zones
(\(\delta w\) = 26.1--34.3 km; bottom left) in our simulations.

\begin{figure}
\centering
\includegraphics[width=1\linewidth,height=\textheight,keepaspectratio,alt={Coupled
kinetic, kinematic, and rheological controls on slab descent velocity
through the 660 (\textbackslash vec\{v\}\_\textbackslash mathrm\{max\},
top row), kinetic and thermal contributions to 660 displacement
(\textbackslash delta z\_k, \textbackslash delta z\_t, second and third
rows), and 660 width (\textbackslash delta w, bottom row) as a function
of kinetic prefactor Z\_\textbackslash mathrm\{ri\} across the SW suite
(left column) and the NC suite (right column) after 100 Ma (Table S1).
Colors and point shapes denote the four combinations of inflow velocity
(\textbackslash vec\{u\}\_\textbackslash mathrm\{in\} = 5 and 10 cm
yr\^{}\{-1\}) and slab strength (B = 5 and 10); solid and dashed curves
denote no lower-mantle viscosity jump (b = 1) and a 50-fold jump (b =
50), respectively. Gray shading in the top row marks the stagnation
regime (\textbackslash vec\{v\}\_\textbackslash mathrm\{max\}
\textbackslash leq 0.5 cm
yr\^{}\{-1\}).}]{../figs/structure/interactions-quantitative.png}
\caption{Coupled kinetic, kinematic, and rheological controls on slab
descent velocity through the 660 (\(\vec{v}_\mathrm{max}\), top row),
kinetic and thermal contributions to 660 displacement (\(\delta z_k\),
\(\delta z_t\), second and third rows), and 660 width (\(\delta w\),
bottom row) as a function of kinetic prefactor \(Z_\mathrm{ri}\) across
the SW suite (left column) and the NC suite (right column) after 100 Ma
(Table S1). Colors and point shapes denote the four combinations of
inflow velocity (\(\vec{u}_\mathrm{in}\) = 5 and 10 cm yr\(^{-1}\)) and
slab strength (\(B\) = 5 and 10); solid and dashed curves denote no
lower-mantle viscosity jump (\(b\) = 1) and a 50-fold jump (\(b\) = 50),
respectively. Gray shading in the top row marks the stagnation regime
(\(\vec{v}_\mathrm{max} \leq\) 0.5 cm
yr\(^{-1}\)).}\label{fig:interactions}
\end{figure}

\clearpage

\section{Discussion}\label{sec:discussion}

\subsection{Uncertainties and Model
Limitations}\label{sec:uncertainties-and-model-limitations}

Distinguishing whether subducting slabs operate in the
nominal-equilibrium or sluggish-stagnation regime is currently limited
by both large experimental uncertainties and natural slab variations
that have not yet been systematically accounted for in kinetics
experiments. The primary quantitative uncertainty for the post-spinel
reaction stems from the activation energy \(E^{\ast}\) for Si
self-diffusion and the kinetic prefactor \(Z_\mathrm{ri}\), which
bundles the growth rate prefactor \(K^{\prime}\) and grain size \(d\).
Monte Carlo regression of the Kubo et al. (2002) experimental dataset
yields \(E^{\ast}\) = 310 \(\pm\) 233 kJ mol\(^{-1}\) and
\(\ln(K^{\prime})\) = -29.8 \(\pm\) 22.1 (2\(\sigma\); Figure S3).
Interpreting this range requires distinguishing between systematic
experimental error and random natural variability. Systematic error
arises from uncertainties in controlling and measuring temperature,
pressure, and transformation progress at extreme multi-anvil conditions,
and defines the wide laboratory confidence envelope within which both SW
and NC suites sit (Figure S4). Random variability, by contrast, reflects
physical differences among subducting slabs that can shift effective
\(Z_\mathrm{ri}\) by several orders of magnitude from one subduction
zone to another. Two physical mechanisms dominate this natural
variability alongside bulk Fe-Mg composition: 1) parent ringwoodite
grain size \(d\), which governs the balance between faster
grain-boundary diffusion and slower intracrystalline Si diffusion (Kubo
et al., 2000; Mao \& Zhong, 2021; Yamazaki et al., 2005), and 2)
structurally bound water, which may enhance Si diffusivity and lamellar
growth rates by orders of magnitude relative to anhydrous ringwoodite
(Chakraborty, 2010). Both factors correlate with subduction history,
which our simulations do not account for. Cold, dry slab cores that
undergo high-overstep transformation at the 410 and 520 can experience
grain-size reduction to 1--100 \(\mu\)m. In contrast, warmer or more
slowly subducting slabs that transform near equilibrium retain coarser
millimeter-scale grains (Riedel \& Karato, 1997). We also neglect major
contributing factors that vary among subduction zones such as slab
hydration, faulting, and sediment input. Our two simulation suites only
bracket physically permitted end-members, demonstrating that new
kinetics experiments addressing both the large laboratory uncertainties
of the Kubo et al. (2002) calibration and natural variations in grain
size, water, and Fe-Mg content in slabs are crucial before kinetic
regimes can be independently correlated with individual subduction
zones.

Another key driver of uncertainty is that the kinetics experiments of
Kubo et al. (2002) were performed at highly overstepped metastable
conditions (\(>\) 1 GPa), requiring significant extrapolation to
near-equilibrium conditions in our simulations. The quadratic dependence
of ringwoodite decomposition on excess Gibbs energy
(Equation~\ref{eq:reaction-rate-660}) results in vanishingly small
reaction rates at small pressure oversteps in our simulations. Testing
whether this quadratic suppression persists near equilibrium or gives
way to an alternative, faster transformation mechanism in natural
systems makes new kinetics experiments at small pressure oversteps a
critical priority.

A related simplification concerns the reaction's directionality.
Although Equation~\ref{eq:reaction-rate-660} preserves the sign of
\(\Delta G\) for bidirectional integration, treating the forward
decomposition (ringwoodite \(\Rightarrow\) bridgmanite + periclase) and
reverse recombination (bridgmanite + periclase \(\Rightarrow\)
ringwoodite) symmetrically is an approximation. Experiments show that
the reverse reaction proceeds via a different mechanism and
substantially more slowly than forward decomposition below 2000 K
(Dobson \& Mariani, 2014; Lessing et al., 2022), introducing a kinetic
asymmetry not captured by our symmetric \(\Delta G \cdot |\Delta G|\)
formulation. Because a symmetric rate law would overestimate reverse
reaction rates in ascending plumes and underestimate metastable
bridgmanite + periclase persistence in upwelling material, we restrict
this study to forward decomposition in subducting slabs and leave
upwelling plumes to future work incorporating dedicated reverse-reaction
kinetics.

This kinetic treatment is itself coupled to thermodynamics in ways we
only partially resolve. While our material model implicitly includes
latent heating near equilibrium, the latent-heating effect at
overstepped metastable conditions merits particular caution for the
post-spinel reaction. Because the ringwoodite \(\Leftrightarrow\)
bridgmanite + periclase Clapeyron slope is negative, the (endothermic)
latent-heat effect close to equilibrium reinforces kinetic inhibition.
At highly overstepped metastable conditions, however, the free energy
from mechanical work \(P\Delta V\) (rather than \(T\Delta S\)) becomes
operative and this reinforcement is lost. This contrasts with the 410,
where both the (exothermic) latent-heat effect and \(P\Delta V\)
contributions promote faster kinetics. Fully resolving this feedback
would likely strengthen the regime contrasts we report here and merits
dedicated investigation. We also define Gibbs free energy based on the
reference pressure (Equation~\ref{eq:gibbs}), neglecting deviatoric
stresses at individual grain interfaces that directly influence reaction
pathways and chemical potentials (Wheeler, 2014, 2020). While the
reference pressure can provide a reasonable approximation under certain
conditions (Wheeler, 2015), an improved approach would be to use the
dynamic pressure, i.e.~the mean normal stress actually experienced by
the rock. A rigorous treatment will require consideration of more
localized stress-dependent chemical potentials (work in progress).

Beyond these kinetic and thermodynamic uncertainties, our compositional
choices introduce additional simplifications. Our pure Mg end-member
composition neglects Fe-Mg solid solutions, akimotoite, garnet, and
other transition-zone phases present in natural systems (Ita \&
Stixrude, 1992; Papanagnou et al., 2023; Xu et al., 2008). Iron is
unlikely to alter our conclusions substantially because Si
self-diffusion limits ringwoodite decomposition (Kubo et al., 2002;
Yamazaki et al., 2005). Iron content should therefore have only a minor
effect on reaction kinetics, and the narrow (\(<\) 1 kbar) Fe-Mg
post-spinel phase loop ({Ishii et al.}, 2019) means metastable
Fe-bearing ringwoodite likely transforms directly to the stable
bridgmanite + periclase assemblage, keeping the transition effectively
univariant as in our pure Mg model. In the coldest slab interiors, the
akimotoite \(\Leftrightarrow\) bridgmanite transition may dominate over
the post-spinel transition ({Chanyshev et al.}, 2022; Cottaar \& Deuss,
2016), producing additional deepening beyond our kinetic model. This
reaction is nevertheless expected to be comparatively rapid because the
ilmenite and perovskite structures are crystallographically closely
related (Navrotsky, 1998).

A final set of simplifications concerns the large-scale viscosity
structure and dimensionality of our simulations. Our simulations test
lower-mantle viscosity jumps of \(b\) = 1 and \(b\) = 50 to partially
characterize how viscosity structure interacts with kinetic inhibition
(Section~\ref{sec:interactions-among-key-physical-parameters}; see
Figures S6--S7 and S13--S14). This approach does not span the full
plausible range of lower-mantle viscosity contrasts (estimated at
roughly 10--100\(\times\), Karato, 2008; {Goes et al.}, 2017), and
quantitative predictions remain model-dependent. More broadly, natural
subduction involves the simultaneous action of viscosity structure, the
Clapeyron slope, plate forcing, and kinetics ({Agrusta et al.}, 2017;
Garel et al., 2014; {Goes et al.}, 2017), and isolating each
contribution requires independent observational constraints. Grain-size
distribution at the post-spinel reaction front (Kubo et al., 2000; Mao
\& Zhong, 2021; Yamazaki et al., 2005) could also generate viscosity
contrasts that augment stagnation beyond what our rheological model
captures. Our simple temperature-dependent rheological model omits
grain-size evolution that affects both flow dynamics and reaction
kinetics (Hirth \& Kohlstedf, 2003; Karato, 1984; Yamazaki et al.,
2005). Finally, our 2D simulations neglect 3D slab geometry, trench
migration, slab rollback, and lateral flow diversion. In 3D, a slab that
retains its dip and is carried laterally away from the 660 by rollback
could sustain faster descent than the vertically ponding slabs in our 2D
domain, since less metastable material would accumulate locally beneath
any one point on the slab. These 3D effects, alongside slab morphology
and residence time more broadly, warrant further investigation (Garel et
al., 2014; {Sime et al.}, 2024). Despite these limitations, our results
capture the first-order kinetic effects governing 660 structure and slab
dynamics across three distinct regimes.

\subsection{Contrasting Reaction Mechanisms at the 410 and
660}\label{sec:contrasting-reaction-mechanisms-at-the-410-and-660}

Unlike interface-controlled 410 kinetics, which produce
threshold-controlled re-sharpening, diffusion-controlled 660 kinetics
drive progressive and persistent slab stagnation across three regimes
(Equations \ref{eq:volume-fraction} to \ref{eq:reaction-rate-660}). At
the 410, the interface-controlled driving-force term
\((1 - \exp[-\Delta G / RT])\) saturates toward unity once
\(|\Delta G| \gg RT\) (Equation~\ref{eq:reaction-rate-410}). Beyond this
saturation point, only the Arrhenius term
\(\exp[-(H^\ast + PV^\ast)/RT]\) controls reaction progress. In a cold,
stagnated slab, this factor remains negligible until temperatures rise
enough to activate transformation, at which point the reaction completes
abruptly over a narrow depth interval, re-sharpening the 410. This
threshold produced the three regimes (quasi-equilibrium, intermediate,
and ultra-sluggish with re-sharpening) with distinct 410 seismic
signatures (Kerswell et al., 2026).

At the 660, the diffusion-controlled driving-force term
\(\Delta G \cdot |\Delta G|\) grows quadratically without bound as
ringwoodite is carried below its stability field. Larger overstepping
always drives larger (silicon) chemical potential gradients across the
decomposition lamellae and thus faster transformation (Dobson \&
Mariani, 2014; Kubo et al., 2002; Yamazaki et al., 2005; Zener, 1946).
No saturation ceiling exists, so the driving force never yields to the
Arrhenius term alone. As a result, \(\dot{X}_\mathrm{max}\) increases
progressively with depth and the 660 widens (\(\delta w\)) and deepens
(\(\delta z_k\)) continuously across all three regimes without
triggering an abrupt transformation. The distinction between the
intermediate and stagnation regimes at the 660 is therefore a matter of
degree, marked by progressively longer reaction timescales \(\tau\)
(Equation~\ref{eq:reaction-timescale}), larger reaction length scales
\(L\) (Equation~\ref{eq:characteristic-length-scale}), slower descent,
and broader, more deeply displaced discontinuities, rather than the
abrupt threshold that separates the ultra-sluggish regime at the 410
(Figure~\ref{fig:contributions}).

This unbounded quadratic driving force also explains why the NC suite
remains in quasi-equilibrium at the 660 across nearly the entire nominal
\(Z_\mathrm{ri}\) range, even when sluggish 410 prefactors in the same
cold slabs (\(\Delta T_0\) = -1000 K) generate deep metastable olivine
wedges (test scenarios 16--19, T1--T4; Table S1). Colder slab cores at
660 km (\(\hat{T} \approx\) -358 to -839 K in NC versus -236 to -457 K
in SW) force a large Arrhenius penalty via \(\exp(-E^{\ast} / RT)\), but
the higher nominal \(Z_\mathrm{ri}\) values compensate so that the
product \(Z_\mathrm{ri}\exp(-E^{\ast} / RT)\) at the fastest-SW and
slowest-NC end-members overlaps (Figure~\ref{fig:contributions}).
Because colder slab cores also experience larger thermal displacement
(\(\delta z_t\) = 16.6--39.0 km) and stronger overstepping
\(|\Delta G| \approx |\hat{T}\, \Delta \bar{S}|\), the quadratic
\(\Delta G \cdot |\Delta G|\) factor accelerates decomposition in the NC
suite by up to five orders of magnitude relative to the slow end of the
SW suite (Figure~\ref{fig:contributions}). Whereas saturation of
\((1 - \exp[-\Delta G / RT])\) at the 410 allows cold temperatures to
freeze metastable olivine, overstepping at the 660 accelerates
decomposition of ringwoodite, so sustained metastable ringwoodite and
kinetic stagnation require sluggish \(Z_\mathrm{ri}\) prefactors near
the lower bound of the Kubo et al. (2002) envelope (SW suite; Figure
S4).

This difference in reaction mechanism also determines long-term slab
behavior. Metastable olivine accumulated in the ultra-sluggish 410
regime eventually transforms in a rapid pulse that can disrupt transient
stalling or ponding (Kerswell et al., 2026). Metastable ringwoodite in
the 660 stagnation regime is consumed gradually (\(\tau\) up to 26 Ma;
Equation~\ref{eq:reaction-timescale}), sustaining persistent buoyancy
that resists penetration over much longer timescales. Kinetically
stagnated slabs at the 660 are therefore more stable than those at the
410, consistent with the long-lived sub-horizontal ponding observed in
the western Pacific (Fukao \& Obayashi, 2013; Tauzin et al., 2017).

Near equilibrium, the two reaction mechanisms also differ. The quadratic
driving force at the 660 strongly suppresses reaction rates near the
phase boundary because \(\Delta G \cdot |\Delta G|\) scales as
\(\Delta G^2\), approaching zero as \(\Delta G \to\) 0. The linear
driving force term governing the 410
\((1 - \exp[-\Delta G / RT]) \approx \Delta G / RT\) suppresses reaction
rates less strongly when \(\Delta G\) approaches zero. This inherent
near-equilibrium suppression produces a finite baseline kinetic
contribution (\(\delta z_k \approx\) 6--11 km) and broadens the 660
transition zone slightly relative to the 410, even under
quasi-equilibrium conditions (Figure~\ref{fig:contributions};
Figure~\ref{fig:seismic-profiles}).

\subsection{Kinetic Amplification of Temperature-Driven 660
Displacement}\label{sec:kinetic-amplification-of-temperature-driven-660-displacement}

The kinetic amplification quantified in
Section~\ref{sec:structure-of-the-660} (Figure~\ref{fig:contributions})
reconciles observed 30--60 km depressions beneath subduction zones with
modest Clapeyron slopes. Within the pure Mg, garnet-free system we
model, equilibrium thermodynamics predicts downward 660 displacement
from two principal sources: 1) the negative Clapeyron slope shifting the
post-spinel reaction to higher pressure beneath cold slabs, and 2) the
akimotoite \(\Leftrightarrow\) bridgmanite transition generating
additional displacement of 30--90 km in the coldest slab interiors
({Chanyshev et al.}, 2022; Cottaar \& Deuss, 2016). Although our
simulations do not prescribe constant Clapeyron slopes, effective slopes
are implicit in the tabulated molar entropy (\(\Delta \bar{S}\)) and
volume (\(\Delta \bar{V}\)) changes along our pure-Mg thermodynamic
profiles (Figure S2). Evaluating
\(\gamma = \Delta \bar{S} / \Delta \bar{V}\) at the nominal equilibrium
crossings (\(\Delta \bar{G}\) = 0) yields implicit Clapeyron slopes of
\(\gamma_{410}\) = +3.09 MPa K\(^{-1}\), \(\gamma_{520}\) = +5.36 MPa
K\(^{-1}\), and \(\gamma_{660}\) = -1.87 MPa K\(^{-1}\), with
\(\gamma_{660}\) shallowing slightly from -1.87 to -1.60 MPa K\(^{-1}\)
across the 0--80 km displacement range below the nominal 660. Our
implicit post-spinel slope of -1.87 MPa K\(^{-1}\) sits between the
recent experimental estimates (-0.4 to -1.3 MPa K\(^{-1}\), {Katsura et
al.}, 2003; Fei et al., 2004; Ishii et al., 2018; Kojitani et al., 2016;
Litasov et al., 2005) and steeper values reported for earlier
calibrations (-2 to -3 MPa K\(^{-1}\), Ito \& Takahashi, 1989; Bina \&
Helffrich, 1994) or hydrous compositions (Ghosh et al., 2013). Because
-1.87 MPa K\(^{-1}\) is about 1.4 to 4.7 times steeper than modest
Clapeyron slopes (-0.4 to -1.3 MPa K\(^{-1}\)), the thermal
contributions in our simulations (\(\delta z_t\) = 11--21 km in SW and
17--39 km in NC) represent conservative upper bounds, where the same
slab-core thermal anomalies would produce only \(\sim\) 2--15 km (SW) or
\(\sim\) 4--27 km (NC) of thermal displacement (cf.
Figure~\ref{fig:contributions}; Figure~\ref{fig:interactions}; Table
S1). Sluggish kinetics in the SW suite add \(\delta z_k\) = 8.7--64.1 km
in the intermediate and stagnation regimes, amplifying total 660
displacement to 27--80 km (Figure~\ref{fig:seismic-profiles}; Table S1).
Such amplifications are roughly two to five times the thermal
contribution \(\delta z_t\) in our pure-Mg material model
(\(\gamma_{660}\) = -1.87 MPa K\(^{-1}\)) and up to an order of
magnitude larger than experimentally-derived predictions
(\(\gamma_{660}\) = -0.4 to -1.3 MPa K\(^{-1}\)). Beyond these thermal
and kinetic effects, two compositional factors outside our model could
further modify the 660: 1) water in ringwoodite broadens the post-spinel
reaction and raises its pressure (Ghosh et al., 2013; Muir et al.,
2021), promoting additional deepening in cold slabs in contrast to
water's shallowing effect at the 410 (Smyth, 1987; Smyth \& Frost, 2002;
{van der Meijde et al.}, 2003), while 2) basalt accumulation modifies
reflectivity without systematically affecting depth ({Goes et al.},
2022; {Yu et al.}, 2023).

The directional contrast between kinetic effects at the 410 and 660 in
cold slabs is fundamental. At the 410, the positive Clapeyron slope
(\(\gamma_{410}\) = +3.09 MPa K\(^{-1}\)) produces upward thermal
displacement under quasi-equilibrium conditions. Decreasing olivine
\(\Leftrightarrow\) wadsleyite reaction rates first broaden and deepen
the 410 discontinuity, then ultimately re-sharpen it at a critical
temperature threshold, reversing its displacement sign relative to the
equilibrium prediction (Kerswell et al., 2026). At the 660, both the
equilibrium thermal contribution \(\delta z_t\) and the kinetic
contribution \(\delta z_k\) across all three regimes produce downward
displacement (Figure~\ref{fig:contributions};
Figure~\ref{fig:seismic-profiles}). Kinetics amplify rather than reverse
the equilibrium thermal signal at the 660, and the transition from one
regime to the next is monotonic. The 660 therefore lacks the seismically
diagnostic sign-reversal that marks the onset of ultra-sluggish
conditions at the 410.

Jointly observing both discontinuities in the same slab column exploits
these different behaviors. A sharp, deeply displaced 410 paired with a
broad, deeply displaced 660 points to ultra-sluggish olivine
\(\Leftrightarrow\) wadsleyite kinetics coexisting with intermediate to
ultra-sluggish ringwoodite decomposition. A broad, deepened 410 paired
with a broad, deepened 660 is consistent with intermediate kinetic
inhibition at both transitions. A modestly displaced, narrow 410 paired
with a modestly displaced, narrow 660 indicates quasi-equilibrium
conditions across the full transition zone, as in the NC suite
(Figure~\ref{fig:seismic-profiles}). These joint patterns exploit the
different saturation behaviors of the interface-controlled and
diffusion-controlled reaction mechanisms, though they only narrow
(rather than eliminate) the ambiguity with thermal, compositional, and
rheological contributions discussed below.

\begin{figure}
\centering
\includegraphics[width=0.62\linewidth,height=\textheight,keepaspectratio,alt={Synthetic
compressional wave velocity (V\_p) profiles extracted along
representative slab depth profiles after 100 Ma in the SW suite (left)
and the NC suite (right) under identical kinematic and rheological
settings (scenarios 4 and 12, respectively; Table~). Black curves show
the initial adiabatic reference profile, and horizontal gray bands mark
the nominal equilibrium 410, 520, and 660 transition depths. Colored
curves span slow (dark red) to fast (yellow) Z\_\textbackslash
mathrm\{ri\} within each suite. In the SW suite, cold temperatures (-T)
uplift the 410 and 520 while sluggish kinetics (-Z\_\textbackslash
mathrm\{ol\}, -Z\_\textbackslash mathrm\{wd\}) counter-shift them
downward. In contrast, both temperature and sluggish kinetics (-T,
-Z\_\textbackslash mathrm\{ri\}) deepen and broaden the 660
(\textbackslash delta z = 25--48 km; \textbackslash delta z\_k =
8.6--37.4 km; \textbackslash delta w = 6.0--13.4 km). In the NC suite,
colder temperatures (-T) strongly uplift the 410 and 520 and deepen the
660 by \textbackslash delta z = 37--40 km (\textbackslash delta z\_t =
29.6 km), while all Z\_\textbackslash mathrm\{ri\} curves collapse onto
a single narrow band (\textbackslash delta z\_k = 7.8--10.5 km;
\textbackslash delta w = 6.3--8.1
km).}]{../figs/seismic/seismic-sw-nc.png}
\caption{Synthetic compressional wave velocity (\(V_p\)) profiles
extracted along representative slab depth profiles after 100 Ma in the
SW suite (left) and the NC suite (right) under identical kinematic and
rheological settings (scenarios 4 and 12, respectively;
Table~\ref{tbl:scenarios}). Black curves show the initial adiabatic
reference profile, and horizontal gray bands mark the nominal
equilibrium 410, 520, and 660 transition depths. Colored curves span
slow (dark red) to fast (yellow) \(Z_\mathrm{ri}\) within each suite. In
the SW suite, cold temperatures (\(-T\)) uplift the 410 and 520 while
sluggish kinetics (\(-Z_\mathrm{ol}\), \(-Z_\mathrm{wd}\)) counter-shift
them downward. In contrast, both temperature and sluggish kinetics
(\(-T\), \(-Z_\mathrm{ri}\)) deepen and broaden the 660 (\(\delta z\) =
25--48 km; \(\delta z_k\) = 8.6--37.4 km; \(\delta w\) = 6.0--13.4 km).
In the NC suite, colder temperatures (\(-T\)) strongly uplift the 410
and 520 and deepen the 660 by \(\delta z\) = 37--40 km (\(\delta z_t\) =
29.6 km), while all \(Z_\mathrm{ri}\) curves collapse onto a single
narrow band (\(\delta z_k\) = 7.8--10.5 km; \(\delta w\) = 6.3--8.1
km).}\label{fig:seismic-profiles}
\end{figure}

\subsection{Constraining Kinetics from Seismic
Observations}\label{sec:constraining-kinetics-from-seismic-observations}

Kinetic inhibition offers a potential explanation for broad, deeply
displaced 660 structures beneath subduction zones that are commonly
attributed to water or the akimotoite \(\Leftrightarrow\) bridgmanite
transition alone. Seismic observations beneath South America (Ba et al.,
2025), East Asia (Liu et al., 2022; Sun et al., 2020; Tang et al.,
2023), and Europe (Cottaar \& Deuss, 2016; Kalmár et al., 2025) show the
660 displaced over a broad area while 410 displacement is either more
localized or undetected, possibly because narrow structures escape
detection.

Beyond overall displacement, transition width (\(\delta w\)) provides an
independent observational test of our kinetic-inhibition hypothesis
because sluggish kinetics alone can reproduce much of the broadening
often attributed to combined ringwoodite and garnet breakdown. Beneath
the western US, receiver-function inversions resolve a \(\sim\) 20 km
wide 410 alongside a sharp 660 with an enhanced velocity gradient
directly above it associated with stalled Farallon fragments (Schmandt,
2012), placing an upper bound on local 660 kinetic broadening. Beneath
Chile--Argentina, frequency-dependent receiver functions constrain a
30--40 km wide 660 (Bonatto et al., 2020), and beneath Northeast China,
triplication modeling resolves a 50--70 km transition (Wang \& Niu,
2010); both locations are associated with subducted slabs. The 30--40 km
width beneath Chile--Argentina overlaps the upper end of the range
predicted by sluggish kinetics in the SW suite (up to \(\delta w\) = 34
km in strong, rapidly descending slabs; Figure~\ref{fig:contributions};
Figure~\ref{fig:interactions}; Table S1). The 50--70 km transition
beneath Northeast China exceeds our pure-Mg kinetic bound, however, and
points to sluggish kinetics acting alongside garnet dissolution or the
akimotoite \(\Leftrightarrow\) bridgmanite transition.

Table~\ref{tbl:regional-observations} tests this thermal-kinetic
interpretation against representative subduction zones, reporting the
conservative reading where methods disagree. Penetrating slabs show
modest 660 displacements consistent with quasi-equilibrium thermal
contributions \(\delta z_t\) for experimentally-determined Clapeyron
slopes ({Goes et al.}, 2017), as in our NC suite. Stagnant slabs
commonly show displacements of 20--60 km with 660 gradients broadened
over tens of kilometers, consistent with intermediate to
stagnation-regime kinetics in the SW suite where \(\delta z_k\)
amplifies modest thermal displacements. Across our simulations, kinetic
effects reach up to \(\delta z\) = 80 km total displacement
(\(\delta z_k\) up to 64 km) and \(\delta w\) = 34 km width in the
strongest, fastest SW slabs (Figure~\ref{fig:interactions}; Table S1).
Signals that exceed those bounds, notably the 50--70 km wide 660 beneath
Northeast China (Wang \& Niu, 2010) and double-step displacements
reaching \(\sim\) 740 km beneath Tonga--Kermadec ({Chanyshev et al.},
2022), therefore require additional contributions from the akimotoite
\(\Leftrightarrow\) bridgmanite transition ({Chanyshev et al.}, 2022;
Cottaar \& Deuss, 2016) or combined ringwoodite and garnet breakdown
(Wang \& Niu, 2010). Water in ringwoodite can also deepen the 660 (Muir
et al., 2021; Sun et al., 2020), although that hypothesis fails
complementary tests beneath central Europe (Cottaar \& Deuss, 2016).

\landscape

\begin{longtable}[]{@{}
  >{\raggedright\arraybackslash}p{(\linewidth - 10\tabcolsep) * \real{0.1635}}
  >{\raggedright\arraybackslash}p{(\linewidth - 10\tabcolsep) * \real{0.2308}}
  >{\raggedright\arraybackslash}p{(\linewidth - 10\tabcolsep) * \real{0.1538}}
  >{\raggedright\arraybackslash}p{(\linewidth - 10\tabcolsep) * \real{0.1058}}
  >{\raggedright\arraybackslash}p{(\linewidth - 10\tabcolsep) * \real{0.2885}}
  >{\raggedright\arraybackslash}p{(\linewidth - 10\tabcolsep) * \real{0.0577}}@{}}
\caption{Representative regional 660 observations. Displacement is
relative to the nominal 660 km global reference. Width is the depth
extent of the 660 velocity gradient where resolved. References: 1: Fukao
\& Obayashi (2013), 2: Tauzin et al. (2017), 3: Waszek et al. (2024), 4:
Wang \& Niu (2010), 5: Tang et al. (2023), 6: Sun et al. (2020), 7:
{Chanyshev et al.} (2022), 8: {Goes et al.} (2017), 9: Ba et al. (2025),
10: Bonatto et al. (2020), 11: Schmandt (2012), 12: Cottaar \& Deuss
(2016), 13: Kalmár et al.
(2025)}\label{tbl:regional-observations}\tabularnewline
\toprule\noalign{}
\begin{minipage}[b]{\linewidth}\raggedright
Subduction zone
\end{minipage} & \begin{minipage}[b]{\linewidth}\raggedright
Slab morphology at 660
\end{minipage} & \begin{minipage}[b]{\linewidth}\raggedright
660 displacement
\end{minipage} & \begin{minipage}[b]{\linewidth}\raggedright
660 width
\end{minipage} & \begin{minipage}[b]{\linewidth}\raggedright
Interpretation in this study
\end{minipage} & \begin{minipage}[b]{\linewidth}\raggedright
Refs
\end{minipage} \\
\midrule\noalign{}
\endfirsthead
\toprule\noalign{}
\begin{minipage}[b]{\linewidth}\raggedright
Subduction zone
\end{minipage} & \begin{minipage}[b]{\linewidth}\raggedright
Slab morphology at 660
\end{minipage} & \begin{minipage}[b]{\linewidth}\raggedright
660 displacement
\end{minipage} & \begin{minipage}[b]{\linewidth}\raggedright
660 width
\end{minipage} & \begin{minipage}[b]{\linewidth}\raggedright
Interpretation in this study
\end{minipage} & \begin{minipage}[b]{\linewidth}\raggedright
Refs
\end{minipage} \\
\midrule\noalign{}
\endhead
\bottomrule\noalign{}
\endlastfoot
Kuril--Izu--Bonin & Stagnant (variable) & Tens of km & & Intermediate to
stagnation-regime kinetics & 1--3 \\
Northeast China & Stagnant & \(\sim\) 20 km & 50--70 km &
Stagnation-regime kinetics in a strong slab and/or akimotoite & 4--8 \\
Mariana & Stagnant remnants & Modest & & Quasi-equilibrium & 1, 8 \\
Tonga--Kermadec & Penetrating & \(\leq\) 80 km & & Akimotoite (double
step) rather than kinetics & 1, 7 \\
Central America & Penetrating & Modest & & Quasi-equilibrium & 1 \\
Chile--Argentina & Stagnant (variable) & \(\sim\) 15--20 km & 30--40 km
& Stagnation-regime kinetics & 9--10 \\
Western US & Stagnant remnants & & Sharp & Sharp 660 limits kinetic
broadening & 11 \\
Mediterranean & Penetrating (variable) & \(\sim\) 10--20 km & &
Intermediate to stagnation-regime kinetics and/or akimotoite & 7,
12--13 \\
\end{longtable}

\endlandscape

Beyond depth and width, the 660 also displays a structural complexity
that generally exceeds the 410. The discontinuity can appear split
(Boyce \& Cottaar, 2021; Deuss et al., 2006) or difficult to resolve in
long-period PP precursor data (Deuss, 2009; Lessing et al., 2015), and
observed amplitude variations in precursors (Waszek et al., 2021; {Yu et
al.}, 2023) and high-frequency scattering (Day \& Deuss, 2013; Wu et
al., 2019) point to compositional variations around the boundary. Some
of this apparent complexity may itself be a kinetic signature. In
higher-frequency data, such as receiver functions, a broad transition
may resemble a split one. Determining how much kinetic effects
contribute to these observations (versus genuine compositional
variation) requires forward modeling of synthetic seismic data (e.g.,
Lessing et al., 2015).

Seismic observations alone, however, cannot fully separate kinetic
effects from thermal, compositional, and rheological contributions. The
boundaries between regimes shift systematically with \(B\), inflow
velocity, and the viscosity jump factor \(b\), so the same 660 structure
is consistent with many kinetic-rheological combinations. As previously
stated, joint analysis of the 410 and 660 in the same slab column
narrows, rather than eliminates, this ambiguity. However, two laboratory
priorities would translate this approach into quantitative observational
constraints. The first is high-pressure experiments under cold-slab and
small-overstep conditions targeting large laboratory uncertainties (see
Section~\ref{sec:uncertainties-and-model-limitations}). The second is
explicit kinetic models for the reverse post-spinel reaction (e.g.,
Lessing et al., 2022).

\subsection{Slab Dynamics and the Diversity of Subduction
Behavior}\label{sec:slab-dynamics-and-the-diversity-of-subduction-behavior}

Sluggish ringwoodite decomposition kinetics independently promote slab
stagnation at the 660 without requiring trench retreat or a prescribed
lower-mantle viscosity jump, establishing reaction kinetics as an
additional, conditional impedance alongside equilibrium controls. Under
nominal ringwoodite kinetics (NC suite), rapid decomposition maintains
equilibrium and prevents metastable buoyancy from accumulating (even in
very cold slabs; \(\Delta T_0\) = -1000 K), so slabs penetrate
continuously unless slowed and folded by a lower-mantle viscosity jump
(Figure~\ref{fig:slab-composition-scen12}; Figures S13--S14 and
S17--S18). Under sluggish ringwoodite kinetics (SW suite), by contrast,
accumulated metastable buoyancy drives the penetration-to-stagnation
progression detailed in Section~\ref{sec:structure-of-the-660}, from
quasi-equilibrium through intermediate flattening to ponding above the
660 (Figure~\ref{fig:slab-composition-scen04};
Figure~\ref{fig:contributions}). These results complement the
equilibrium stagnation mechanisms identified by Garel et al. (2014),
{Agrusta et al.} (2017), and {Goes et al.} (2017). Where those studies
identify the Clapeyron slope, lower-mantle viscosity increase, and
trench retreat as necessary joint conditions for stagnation, our SW
suite establishes sluggish kinetic inhibition as an additional,
independently variable impedance. A lower-mantle viscosity jump
reinforces kinetic stagnation in our simulations
(Figure~\ref{fig:interactions}; Figures S6--S7 and S10--S11), confirming
that kinetic and viscous impedances at the 660 are additive.

The diversity of slab behavior documented globally by Fukao \& Obayashi
(2013) maps onto this interpretation. Stagnant western Pacific slabs
ponding above the 660 with displacements of 30--60 km (Fukao \&
Obayashi, 2013; Tauzin et al., 2017) are consistent with intermediate to
stagnation-regime kinetics in the SW suite, where kinetic contributions
(\(\delta z_k\) = 9--64 km) amplify modest thermal displacements.
Steeply penetrating slabs with narrow, modestly displaced 660s are
consistent with the quasi-equilibrium regime of both suites. Our results
also help explain slabs that flatten or stagnate in the uppermost lower
mantle between 660 and 1000 km depth, such as segments of the Mariana,
Kermadec, Java, and Central America subduction zones (Fukao \& Obayashi,
2013; {Goes et al.}, 2017). Two complementary mechanisms may play a
role: 1) in strong, rapidly subducting SW slabs under sluggish kinetics
(scenarios 3 and 7; Figures S10--S11), structural coherence resists
buckling at the 660 and advects buoyant metastable ringwoodite to
\(\sim\) 740 km depth before the slab decelerates and flattens in the
uppermost lower mantle; and 2) because metastable ringwoodite decomposes
gradually over characteristic timescales of \(\tau\) \(\sim\) 5--26 Ma
rather than in an abrupt pulse, slabs that initially stall near the base
of the transition zone slowly lose positive buoyancy over 20--160 Ma
(Chapman \& Clarke, 2021; Mao \& Zhong, 2018) and sink into the
uppermost lower mantle, where a 10--100\(\times\) viscosity increase
(scenarios 0--3 and 8--11; Figures S6--S7, S10--S11, S13--S14, and
S17--S18) folds and stalls the slab between 660 and 1000 km depth
(analogous to Garel et al., 2014; {Agrusta et al.}, 2017; {Goes et al.},
2017).

Descent rates offer a further point of comparison with nature, though a
more demanding one given the inherent limitations of our
phenomenological model. In the SW suite, descent rates through the 660
in the quasi-equilibrium regime (0.5--3.2 cm yr\(^{-1}\)) fall below or
at the lower end of the natural interquartile range of 2.6--4.8 cm
yr\(^{-1}\) (Lallemand et al., 2005), a discrepancy also present in the
410 results of Kerswell et al. (2026). The colder, more negatively
buoyant NC suite slabs at \(b\) = 1 (2.9--6.6 cm yr\(^{-1}\)) span the
natural range and slow to 0.9--1.9 cm yr\(^{-1}\) at \(b\) = 50
(Figure~\ref{fig:interactions}; Table S1). Our model includes slab pull,
the dominant driving force in nature, but omits other controls on
descent rate, including upper mantle viscosity structure, chemical
buoyancy contrasts between basaltic and harzburgitic layers, volatile
content, and 2D geometric simplifications that neglect trench retreat
and lateral flow. Any of these factors could shift simulated rates
closer to natural values. Because our simulations are designed to
isolate and demonstrate kinetic effects rather than to reproduce the
full boundary conditions of natural subduction, the absolute descent
rates are not expected to match observations quantitatively. The
qualitative correspondence with observed slab diversity discussed above
is the more direct test of our kinetic-inhibition hypothesis.
Incorporating the omitted physics would be needed to translate our
results into quantitative natural descent rates, and represents a
natural extension of this approach.

Kinetic inhibition acts jointly with slab rheology and kinematics to
control subduction dynamics and 660 structure
(Figure~\ref{fig:interactions}). In the SW suite, weaker slabs
decelerate more readily under metastable buoyancy and stagnate at up to
two orders of magnitude higher \(Z_\mathrm{ri}\) than stronger slabs,
which caps their maximum kinetic contribution (\(\delta z_k\) = 37--59
km) once ponding halts downward advection. Stronger slabs resist
stagnation to lower \(Z_\mathrm{ri}\) and advect untransformed
ringwoodite deeper into the metastable zone when their descent slows,
generating the widest (\(\delta w\) = 26--34 km) and most deeply
displaced (\(\delta z\) = 78--80 km; \(\delta z_k\) = 61--64 km) 660s in
our simulations (Figure~\ref{fig:interactions}; Figures S10--S11). In
the NC suite, conversely, only the weakest and kinematically fastest
configuration (scenario 14; Figures S15--S16) descends rapidly enough at
the slowest nominal \(Z_\mathrm{ri}\) to enter the intermediate regime
(\(\delta z_k\) = 16--22 km). Small changes in convergence rate, slab
thermal state, or rheological strength near a regime boundary can
therefore drive episodic transitions between penetrating and stagnant
behavior at individual subduction zones over geological time ({Agrusta
et al.}, 2017; Kerswell et al., 2026; Mao \& Zhong, 2018).

\section{Conclusions}\label{sec:conclusions}

Diffusion-controlled ringwoodite decomposition kinetics provide a
credible physical mechanism to independently stall subducting slabs and
reconcile deep, broad 660 km discontinuity transitions with modest
experimental Clapeyron slopes. By coupling experimentally calibrated
kinetics to 2D compressible subduction simulations, and decomposing
total 660 displacement into its thermal and kinetic contributions, we
demonstrate that reaction kinetics is an indispensable quantitative
paradigm for interpreting the diversity of observed slab morphologies
and seismic structures across the MTZ.

Ringwoodite decomposition kinetics exert a conditional, first-order
control on slab dynamics and 660 structure, modulated by lower-mantle
viscosity jump, inflow velocity, and slab strength. Nominal kinetic
prefactors (\(Z_\mathrm{ri}\) best-fit to experimental data; NC suite)
create conditions that are indistinguishable from thermodynamic
equilibrium. Slabs penetrate continuously unless slowed by a
lower-mantle viscosity jump, and observed 660 displacements must be
mostly thermal and/or compositional. By contrast, sluggish kinetic
prefactors (representative of dry, coarse-grained rocks; SW suite)
transition slabs from penetration to stagnation with decreasing
\(Z_\mathrm{ri}\). Accumulating metastable ringwoodite slows descent up
to tenfold and ponds slabs above the lower mantle. In these regimes,
kinetic inhibition widens the 660 to 34 km and contributes up to 64 km
of downward displacement (80 km total), amplifying thermal displacements
by two to five times in our pure-Mg model (\(\gamma_{660}\) = -1.87 MPa
K\(^{-1}\)) and by up to an order of magnitude relative to
experimentally-determined Clapeyron slopes (\(\gamma_{660}\) = -0.4 to
-1.3 MPa K\(^{-1}\)). Sluggish kinetics thus act as an independent
stagnation mechanism alongside thermal effects and lower-mantle
viscosity jump, reconciling broad, 30--60 km 660 displacements with
modest Clapeyron slopes without requiring akimotoite, garnet, or water.

Kinetic stagnation at the 660 is progressive and persistent (\(\tau\) up
to 26 Ma) rather than threshold-controlled, a distinction essential for
interpreting the diverse seismic signatures and slab morphologies across
the MTZ. At the 410, the interface-controlled driving force saturates at
high overstepping, allowing cold temperatures to freeze metastable
olivine until warming triggers an abrupt re-sharpening pulse (Kerswell
et al., 2026). At the 660, the quadratic diffusion-controlled driving
force (\(\Delta G \cdot |\Delta G|\)) grows without bound as
overstepping increases, causing kinetic inhibition to progress
monotonically from quasi-equilibrium through stagnation with no
re-sharpening at any tested condition.

Distinguishing between equilibrium and kinetically controlled 660
regimes in Earth's mantle requires narrowing large laboratory
uncertainties alongside new observational constraints. High-pressure
experiments targeting cold-slab conditions and small pressure oversteps
(accounting for expected grain-size, water, and Fe-Mg variations in the
mantle), dedicated kinetic models for the reverse post-spinel reaction,
and frequency-dependent inversions jointly constraining 410 and 660
topography and gradient width, would discriminate kinetic inhibition
from competing thermal and compositional explanations. Until these
fundamental constraints narrow the permitted kinetic range, both NC and
SW suite end-member behaviors remain plausible.

\section*{Acknowledgements}\label{sec:acknowledgements}
\addcontentsline{toc}{section}{Acknowledgements}

This work was funded by the UKRI NERC Large Grant no. NE/V018477/1. All
computations were undertaken on Barkla2, part of the High Performance
Computing facilities at the University of Liverpool, who graciously
provided expert support. We thank the Computational Infrastructure for
Geodynamics (\url{https://geodynamics.org}) which is funded by the
National Science Foundation under awards EAR-0949446 and EAR-1550901 for
supporting the development of ASPECT. We thank two anonymous reviewers
for their feedback, which greatly improved this work, and the JGR:SE
editorial staff for handling the manuscript.

\section*{Data Availability}\label{sec:data-availability}
\addcontentsline{toc}{section}{Data Availability}

All simulation outputs and post-processed data for reproducing this work
are archived on Zenodo (Kerswell, 2027a), as well as all analysis code
(Kerswell, 2027b). All data and code within these repositories are MIT
Licensed and free for use and distribution (see license details). ASPECT
version 3.1.0 (Bangerth et al., 2026) was used for the computations in
this study and is freely available under the GPL v2.0 or later license.

\section*{Conflict of Interest}\label{sec:conflict-of-interest}
\addcontentsline{toc}{section}{Conflict of Interest}

The authors declare there are no conflicts of interest for this
manuscript.

\clearpage

\section*{References}\label{sec:references}
\addcontentsline{toc}{section}{References}

\protect\phantomsection\label{refs}
\begin{CSLReferences}{1}{0}
\bibitem[\citeproctext]{ref-agrusta2017}
{Agrusta, R., Goes, S., \& van Hunen, J.} (2017). Subducting-slab
transition-zone interaction: Stagnation, penetration and mode switches.
\emph{Earth and Planetary Science Letters}, \emph{464}, 10--23.

\bibitem[\citeproctext]{ref-ba2025}
Ba, K., Huang, J., Yu, Y., Song, J., Li, S., Shen, C., \& Liu, L.
(2025). Break-off and stagnation of the subducting {Nazca} slab in the
mantle transition zone beneath {West-Central South America}.
\emph{Tectonophysics}, \emph{911}, 230834.

\bibitem[\citeproctext]{ref-aspectmanual}
Bangerth, W., Dannberg, J., Fraters, M., Gassmöller, R., Glerum, A.,
Heister, T., et al. (2018, May). {ASPECT: Advanced Solver for Planetary
Evolution, Convection, and Tectonics {[}User Manual{]}} (Version 12).
Figshare. \url{https://doi.org/10.6084/m9.figshare.4865333}

\bibitem[\citeproctext]{ref-aspect-doi-v3.1.0}
Bangerth, W., Dannberg, J., Douglas, D., Fraters, M., Gassmöller, R.,
Glerum, A., et al. (2026, September). {ASPECT: Advanced Solver for
Planetary Evolution, Convection, and Tectonics {[}Software{]}} (Version
3.1.0). Zenodo. \url{https://doi.org/10.5281/zenodo.22814863}

\bibitem[\citeproctext]{ref-bina1994}
Bina, C., \& Helffrich, G. (1994). Phase transition {Clapeyron} slopes
and transition zone seismic discontinuity topography. \emph{Journal of
Geophysical Research: Solid Earth}, \emph{99}(B8), 15853--15860.

\bibitem[\citeproctext]{ref-bonatto2020}
Bonatto, L., Piromallo, C., \& Cottaar, S. (2020). The transition zone
beneath {West Argentina}-{Central Chile} using {P-to-S} converted waves.
\emph{Journal of Geophysical Research: Solid Earth}, \emph{125}(8),
e2020JB019446.

\bibitem[\citeproctext]{ref-boyce2021}
Boyce, A., \& Cottaar, S. (2021). Insights into deep mantle
thermochemical contributions to {African} magmatism from converted
seismic phases. \emph{Geochemistry, Geophysics, Geosystems},
\emph{22}(3), e2020GC009478.

\bibitem[\citeproctext]{ref-burke1965}
Burke, J. (1965). The kinetics of phase transformations in metals.

\bibitem[\citeproctext]{ref-cahn1956}
Cahn, J. (1956). The kinetics of grain boundary nucleated reactions.
\emph{Acta Metallurgica}, \emph{4}(5), 449--459.

\bibitem[\citeproctext]{ref-chakraborty2010}
Chakraborty, S. (2010). Diffusion coefficients in olivine, wadsleyite
and ringwoodite. \emph{Reviews in Mineralogy and Geochemistry},
\emph{72}(1), 603--639.

\bibitem[\citeproctext]{ref-chanyshev2022}
{Chanyshev, A., Ishii, T., Bondar, D., Bhat, S., Kim, E., Farla, R., et
al.} (2022). Depressed 660-km discontinuity caused by
akimotoite--bridgmanite transition. \emph{Nature}, \emph{601}(7891),
69--73.

\bibitem[\citeproctext]{ref-chapman2021}
Chapman, T., \& Clarke, G. (2021). Delayed growth of ferropericlase and
bridgmanite controls slab residence at the 660-km discontinuity.
\emph{Journal of Geophysical Research: Solid Earth}, \emph{126}(6),
e2020JB021487.

\bibitem[\citeproctext]{ref-christian2007}
Christian, J. (2007). \emph{The theory of transformations in metals and
alloys}. Elsevier.

\bibitem[\citeproctext]{ref-clevenger2021}
Clevenger, T., \& Heister, T. (2021). Comparison between algebraic and
matrix-free geometric multigrid for a {Stokes} problem on adaptive
meshes with variable viscosity. \emph{Numerical Linear Algebra with
Applications}, \emph{28}(5), e2375.

\bibitem[\citeproctext]{ref-connolly2009}
Connolly, J. (2009). The geodynamic equation of state: What and how.
\emph{Geochemistry, Geophysics, Geosystems}, \emph{10}(10).

\bibitem[\citeproctext]{ref-cottaar2016}
Cottaar, S., \& Deuss, A. (2016). Large-scale mantle discontinuity
topography beneath {Europe}: Signature of akimotoite in subducting
slabs. \emph{Journal of Geophysical Research: Solid Earth},
\emph{121}(1), 279--292.

\bibitem[\citeproctext]{ref-cottaar2014}
Cottaar, S., Heister, T., Rose, I., \& Unterborn, C. (2014). {BurnMan}:
A lower mantle mineral physics toolkit. \emph{Geochemistry, Geophysics,
Geosystems}, \emph{15}(4), 1164--1179.

\bibitem[\citeproctext]{ref-dassler1996}
Däßler, R., Yuen, D., Karato, S., \& Riedel, M. (1996). Two-dimensional
thermo-kinetic model for the olivine--spinel phase transition in
subducting slabs. \emph{Physics of the Earth and Planetary Interiors},
\emph{94}(3-4), 217--239.

\bibitem[\citeproctext]{ref-day2013}
Day, E., \& Deuss, A. (2013). Reconciling {PP} and
{\(P^\prime P^\prime\)} precursor observations of a complex 660 km
seismic discontinuity. \emph{Geophysical Journal International},
\emph{194}(2), 834--838.

\bibitem[\citeproctext]{ref-deuss2009}
Deuss, A. (2009). Global observations of mantle discontinuities using
{SS} and {PP} precursors. \emph{Surveys in Geophysics}, \emph{30}(4),
301--326.

\bibitem[\citeproctext]{ref-deuss2006}
Deuss, A., Redfern, S., Chambers, K., \& Woodhouse, J. (2006). The
nature of the 660-kilometer discontinuity in {Earth}'s mantle from
global seismic observations of {PP} precursors. \emph{Science},
\emph{311}(5758), 198--201.

\bibitem[\citeproctext]{ref-devaux1997}
Devaux, J., Schubert, G., \& Anderson, C. (1997). Formation of a
metastable olivine wedge in a descending slab. \emph{Journal of
Geophysical Research: Solid Earth}, \emph{102}(B11), 24627--24637.

\bibitem[\citeproctext]{ref-dobson2014}
Dobson, D., \& Mariani, E. (2014). The kinetics of the reaction of
majorite plus ferropericlase to ringwoodite: Implications for mantle
upwellings crossing the 660 km discontinuity. \emph{Earth and Planetary
Science Letters}, \emph{408}, 110--118.

\bibitem[\citeproctext]{ref-fei2004}
Fei, Y., Van Orman, J., Li, J., Van Westrenen, W., Sanloup, C., Minarik,
W., et al. (2004). Experimentally determined postspinel transformation
boundary in Mg2SiO4 using MgO as an internal pressure standard and its
geophysical implications. \emph{Journal of Geophysical Research: Solid
Earth}, \emph{109}(B2).

\bibitem[\citeproctext]{ref-flanagan1998}
Flanagan, M., \& Shearer, P. (1998). Global mapping of topography on
transition zone velocity discontinuities by stacking {SS} precursors.
\emph{Journal of Geophysical Research: Solid Earth}, \emph{103}(B2),
2673--2692.

\bibitem[\citeproctext]{ref-fukao2013}
Fukao, Y., \& Obayashi, M. (2013). Subducted slabs stagnant above,
penetrating through, and trapped below the 660 km discontinuity.
\emph{Journal of Geophysical Research: Solid Earth}, \emph{118}(11),
5920--5938.

\bibitem[\citeproctext]{ref-garel2014}
Garel, F., Goes, S., Davies, D., Davies, J., Kramer, S., \& Wilson, C.
(2014). Interaction of subducted slabs with the mantle transition-zone:
A regime diagram from 2-d thermo-mechanical models with a mobile trench
and an overriding plate. \emph{Geochemistry, Geophysics, Geosystems},
\emph{15}(5), 1739--1765.

\bibitem[\citeproctext]{ref-gassmoller2018}
Gassmöller, R., Lokavarapu, H., Heien, E., Puckett, E., \& Bangerth, W.
(2018). Flexible and scalable particle-in-cell methods with adaptive
mesh refinement for geodynamic computations. \emph{Geochemistry,
Geophysics, Geosystems}, \emph{19}(9), 3596--3604.

\bibitem[\citeproctext]{ref-gassmoller2020}
Gassmöller, R., Dannberg, J., Bangerth, W., Heister, T., \& Myhill, R.
(2020). On formulations of compressible mantle convection.
\emph{Geophysical Journal International}, \emph{221}(2), 1264--1280.

\bibitem[\citeproctext]{ref-ghosh2013}
Ghosh, S., Ohtani, E., Litasov, K., Suzuki, A., Dobson, D., \&
Funakoshi, K. (2013). Effect of water in depleted mantle on post-spinel
transition and implication for 660 km seismic discontinuity. \emph{Earth
and Planetary Science Letters}, \emph{371}, 103--111.

\bibitem[\citeproctext]{ref-goes2017}
{Goes, S., Agrusta, R., van Hunen, J., \& Garel, F.} (2017).
Subduction-transition zone interaction: A review. \emph{Geosphere},
\emph{13}(3), 644--664.

\bibitem[\citeproctext]{ref-goes2022}
{Goes, S., Yu, C., Ballmer, M., Yan, J., \& van der Hilst, R.} (2022).
Compositional heterogeneity in the mantle transition zone. \emph{Nature
Reviews Earth \& Environment}, \emph{3}(8), 533--550.

\bibitem[\citeproctext]{ref-green1979}
Green, D., Jaques, L., \& Hibberson, W. (1979). Petrogenesis of
mid-ocean ridge basalts. In \emph{The earth: Its origin, structure and
evolution} (pp. 265--300). Academic Press.

\bibitem[\citeproctext]{ref-heister2017}
Heister, T., Dannberg, J., Gassmöller, R., \& Bangerth, W. (2017). High
accuracy mantle convection simulation through modern numerical
methods--{II}: Realistic models and problems. \emph{Geophysical Journal
International}, \emph{210}(2), 833--851.

\bibitem[\citeproctext]{ref-hillert1957}
Hillert, M. (1957). Pressure induced diffusion and deformation during
precipitation, especially graphitization. \emph{Jernkont. Ann},
\emph{141}, 157--181.

\bibitem[\citeproctext]{ref-hillert1969}
Hillert, M. (1969). The role of interfaces in phase transformations.
\emph{The Mechanism of Phase Transformations in Crystalline Solids},
231--247.

\bibitem[\citeproctext]{ref-hindmarsh2005}
Hindmarsh, A., Brown, P., Grant, K., Lee, S., Serban, R., Shumaker, D.,
\& Woodward, C. (2005). {SUNDIALS}: Suite of nonlinear and
differential/algebraic equation solvers. \emph{ACM Transactions on
Mathematical Software (TOMS)}, \emph{31}(3), 363--396.

\bibitem[\citeproctext]{ref-hirth2003}
Hirth, G., \& Kohlstedf, D. (2003). Rheology of the upper mantle and the
mantle wedge: A view from the experimentalists. \emph{Geophysical
Monograph-American Geophysical Union}, \emph{138}, 83--106.

\bibitem[\citeproctext]{ref-hosoya2005}
Hosoya, T., Kubo, T., Ohtani, E., Sano, A., \& Funakoshi, K. (2005).
Water controls the fields of metastable olivine in cold subducting
slabs. \emph{Geophysical Research Letters}, \emph{32}(17).

\bibitem[\citeproctext]{ref-ishii2018}
Ishii, T., Kojitani, H., \& Akaogi, M. (2018). Phase relations and
mineral chemistry in pyrolitic mantle at 1600--2200 c under pressures up
to the uppermost lower mantle: Phase transitions around the 660-km
discontinuity and dynamics of upwelling hot plumes. \emph{Physics of the
Earth and Planetary Interiors}, \emph{274}, 127--137.

\bibitem[\citeproctext]{ref-ishii2019}
{Ishii, T., Huang, R., Myhill, R., Fei, H., Koemets, I., Liu, Z., et
al.} (2019). Sharp 660-km discontinuity controlled by extremely narrow
binary post-spinel transition. \emph{Nature Geoscience}, \emph{12}(10),
869--872.

\bibitem[\citeproctext]{ref-ita1992}
Ita, J., \& Stixrude, L. (1992). Petrology, elasticity, and composition
of the mantle transition zone. \emph{Journal of Geophysical Research:
Solid Earth}, \emph{97}(B5), 6849--6866.

\bibitem[\citeproctext]{ref-ito1989}
Ito, E., \& Takahashi, E. (1989). Postspinel transformations in the
system Mg2SiO4-Fe2SiO4 and some geophysical implications. \emph{Journal
of Geophysical Research: Solid Earth}, \emph{94}(B8), 10637--10646.

\bibitem[\citeproctext]{ref-ji2004}
Ji, S. (2004). A generalized mixture rule for estimating the viscosity
of solid-liquid suspensions and mechanical properties of polyphase rocks
and composite materials. \emph{Journal of Geophysical Research: Solid
Earth}, \emph{109}(B10).

\bibitem[\citeproctext]{ref-kalmar2025}
Kalmár, D., Petrescu, L., Hetényi, G., Michailos, K., Süle, B., Neagoe,
C., \& Bokelmann, G. (2025). Mantle transition zone analysis using
{P-to-S} receiver functions in the {Alpine--Carpathian--Dinarides}
region: Impact of plumes and slabs. \emph{Geophysical Journal
International}, \emph{243}(1), ggaf313.

\bibitem[\citeproctext]{ref-karato1984}
Karato, S. (1984). Grain-size distribution and rheology of the upper
mantle. \emph{Tectonophysics}, \emph{104}(1-2), 155--176.

\bibitem[\citeproctext]{ref-karato2008}
Karato, S. (2008). Deformation of earth materials. \emph{An Introduction
to the Rheology of Solid Earth}, \emph{463}.

\bibitem[\citeproctext]{ref-katsura2003}
{Katsura, T., Yamada, H., Shinmei, T., Kubo, A., Ono, S., Kanzaki, M.,
et al.} (2003). Post-spinel transition in Mg2SiO4 determined by high
p--t in situ x-ray diffractometry. \emph{Physics of the Earth and
Planetary Interiors}, \emph{136}(1-2), 11--24.

\bibitem[\citeproctext]{ref-kerswell2027}
Kerswell, B. (2027a, January).
{buchanankerswell/kerswell\_et\_al\_660\_kinetics: Publication Release
{[}Dataset{]}} (Version 1.0.0). Zenodo.
\url{https://doi.org/10.5281/zenodo.24120946}

\bibitem[\citeproctext]{ref-kerswell2027b}
Kerswell, B. (2027b, January).
{buchanankerswell/kerswell\_et\_al\_660\_kinetics: Publication Release
{[}Software{]}} (Version 1.0.0). Zenodo.
\url{https://doi.org/10.5281/zenodo.21810740}

\bibitem[\citeproctext]{ref-kerswell2026}
Kerswell, B., Wheeler, J., Gassmöller, R., Davies, J., Papanagnou, I.,
\& Cottaar, S. (2026). Beyond equilibrium: Kinetic thresholds and
rheological feedbacks create a potentially complex 410 in slab regions.
\emph{Journal of Geophysical Research: Solid Earth}, \emph{131}(7),
e2026JB033781.

\bibitem[\citeproctext]{ref-kojitani2016}
Kojitani, H., Inoue, T., \& Akaogi, M. (2016). Precise measurements of
enthalpy of postspinel transition in Mg2SiO4 and application to the
phase boundary calculation. \emph{Journal of Geophysical Research: Solid
Earth}, \emph{121}(2), 729--742.

\bibitem[\citeproctext]{ref-kronbichler2012}
Kronbichler, M., Heister, T., \& Bangerth, W. (2012). High accuracy
mantle convection simulation through modern numerical methods.
\emph{Geophysical Journal International}, \emph{191}(1), 12--29.

\bibitem[\citeproctext]{ref-kubo2000}
Kubo, T., Ohtani, E., Kato, T., Urakawa, S., Suzuki, A., Kanbe, Y., et
al. (2000). Formation of metastable assemblages and mechanisms of the
grain-size reduction in the postspinel transformation of
{Mg\(_2\)SiO\(_4\)}. \emph{Geophysical Research Letters}, \emph{27}(6),
807--810.

\bibitem[\citeproctext]{ref-kubo2002}
Kubo, T., Ohtani, E., Kato, T., Urakawa, S., Suzuki, A., Kanbe, Y., et
al. (2002). Mechanisms and kinetics of the post-spinel transformation in
{Mg\(_2\)SiO\(_4\)}. \emph{Physics of the Earth and Planetary
Interiors}, \emph{129}(1-2), 153--171.

\bibitem[\citeproctext]{ref-lallemand2005}
Lallemand, S., Heuret, A., \& Boutelier, D. (2005). On the relationships
between slab dip, back-arc stress, upper plate absolute motion, and
crustal nature in subduction zones. \emph{Geochemistry, Geophysics,
Geosystems}, \emph{6}(9).

\bibitem[\citeproctext]{ref-lebedev2002}
Lebedev, S., Chevrot, S., \& Hilst, R. van der. (2002). Seismic evidence
for olivine phase changes at the 410-and 660-kilometer discontinuities.
\emph{Science}, \emph{296}(5571), 1300--1302.

\bibitem[\citeproctext]{ref-lessing2015}
Lessing, S., Thomas, C., Saki, M., Schmerr, N., \& Vanacore, E. (2015).
On the difficulties of detecting {PP} precursors. \emph{Geophysical
Journal International}, \emph{201}(3), 1666--1681.

\bibitem[\citeproctext]{ref-lessing2022}
Lessing, S., Dobson, D., Rost, S., Cobden, L., \& Thomas, C. (2022).
Kinetic effects on the 660-km-phase transition in mantle upstreams and
seismological implications. \emph{Geophysical Journal International},
\emph{231}(2), 877--893.

\bibitem[\citeproctext]{ref-li2026}
Li, H., \& Billen, M. (2026). A unified framework for slab dynamics:
Coupled metastable kinetics and rheological evolution in subduction
zones. \emph{Physics of the Earth and Planetary Interiors}, 107614.

\bibitem[\citeproctext]{ref-litasov2005}
Litasov, K., Ohtani, E., Sano, A., Suzuki, A., \& Funakoshi, K. (2005).
In situ x-ray diffraction study of post-spinel transformation in a
peridotite mantle: Implication for the 660-km discontinuity. \emph{Earth
and Planetary Science Letters}, \emph{238}(3-4), 311--328.

\bibitem[\citeproctext]{ref-liu2022}
Liu, L., Gao, S., Liu, K., Griffin, W., Li, S., Tong, S., \& Ning, J.
(2022). Mantle dynamics of the {North China Craton}: New insights from
mantle transition zone imaging constrained by {P-to-S} receiver
functions. \emph{Geophysical Journal International}, \emph{231}(1),
629--637.

\bibitem[\citeproctext]{ref-mao2018}
Mao, W., \& Zhong, S. (2018). Slab stagnation due to a reduced viscosity
layer beneath the mantle transition zone. \emph{Nature Geoscience},
\emph{11}(11), 876--881.

\bibitem[\citeproctext]{ref-mao2021}
Mao, W., \& Zhong, S. (2021). Formation of horizontally deflected slabs
in the mantle transition zone caused by spinel-to-post-spinel phase
transition, its associated grainsize reduction effects, and trench
retreat. \emph{Geophysical Research Letters}, \emph{48}(15),
e2021GL093679.

\bibitem[\citeproctext]{ref-muir2021}
Muir, J., Zhang, F., \& Brodholt, J. (2021). The effect of water on the
post-spinel transition and evidence for extreme water contents at the
bottom of the transition zone. \emph{Earth and Planetary Science
Letters}, \emph{565}, 116909.

\bibitem[\citeproctext]{ref-myhill2023}
Myhill, R., Cottaar, S., Heister, T., Rose, I., Unterborn, C., Dannberg,
J., \& Gassmöller, R. (2023). {BurnMan}--a {Python} toolkit for
planetary geophysics, geochemistry and thermodynamics.

\bibitem[\citeproctext]{ref-navrotsky1998}
Navrotsky, A. (1998). Energetics and crystal chemical systematics among
ilmenite, lithium niobate, and perovskite structures. \emph{Chemistry of
Materials}, \emph{10}(10), 2787--2793.

\bibitem[\citeproctext]{ref-papanagnou2023}
Papanagnou, I., Schuberth, B., \& Thomas, C. (2023). Geodynamic
predictions of seismic structure and discontinuity topography of the
mantle transition zone. \emph{Geophysical Journal International},
\emph{234}(1), 355--378.

\bibitem[\citeproctext]{ref-porter2009}
Porter, D., \& Easterling, K. (2009). \emph{Phase transformations in
metals and alloys}. CRC press.

\bibitem[\citeproctext]{ref-reynolds2023}
Reynolds, D., Gardner, D., Woodward, C., \& Chinomona, R. (2023).
{ARKODE}: A flexible {IVP} solver infrastructure for one-step methods.
\emph{ACM Transactions on Mathematical Software}, \emph{49}(2), 1--26.

\bibitem[\citeproctext]{ref-riedel1997}
Riedel, M., \& Karato, S. (1997). Grain-size evolution in subducted
oceanic lithosphere associated with the olivine--spinel transformation
and its effects on rheology. \emph{Earth and Planetary Science Letters},
\emph{148}(1-2), 27--43.

\bibitem[\citeproctext]{ref-ringwood1991}
Ringwood, A. (1991). Phase transformations and their bearing on the
constitution and dynamics of the mantle. \emph{Geochimica et
Cosmochimica Acta}, \emph{55}(8), 2083--2110.

\bibitem[\citeproctext]{ref-rubie1994}
Rubie, D., \& Ross, C. (1994). Kinetics of the olivine--spinel
transformation in subducting lithosphere: Experimental constraints and
implications for deep slab processes. \emph{Physics of the Earth and
Planetary Interiors}, \emph{86}(1-3), 223--243.

\bibitem[\citeproctext]{ref-schmandt2012}
Schmandt, B. (2012). Mantle transition zone shear velocity gradients
beneath {USArray}. \emph{Earth and Planetary Science Letters},
\emph{355}, 119--130.

\bibitem[\citeproctext]{ref-sime2024}
{Sime, N., Wilson, C., \& van Keken, P.} (2024). Thermal modeling of
subduction zones with prescribed and evolving 2D and 3D slab geometries.
\emph{Progress in Earth and Planetary Science}, \emph{11}(1), 14.

\bibitem[\citeproctext]{ref-smyth1987}
Smyth, J. (1987). {\(\beta\)-Mg\(_2\)SiO\(_4\)}; a potential host for
water in the mantle? \emph{American Mineralogist}, \emph{72}(11-12),
1051--1055.

\bibitem[\citeproctext]{ref-smyth2002}
Smyth, J., \& Frost, D. (2002). The effect of water on the 410-km
discontinuity: An experimental study. \emph{Geophysical Research
Letters}, \emph{29}(10), 123--1.

\bibitem[\citeproctext]{ref-stixrude2022}
Stixrude, L., \& Lithgow-Bertelloni, C. (2022). Thermal expansivity,
heat capacity and bulk modulus of the mantle. \emph{Geophysical Journal
International}, \emph{228}(2), 1119--1149.

\bibitem[\citeproctext]{ref-sun2020}
Sun, M., Gao, S., Liu, K., \& Fu, X. (2020). Upper mantle and mantle
transition zone thermal and water content anomalies beneath {NE Asia}:
Constraints from receiver function imaging of the 410 and 660 km
discontinuities. \emph{Earth and Planetary Science Letters}, \emph{532},
116040.

\bibitem[\citeproctext]{ref-tang2014}
Tang, Y., Obayashi, M., Niu, F., Grand, S., Chen, Y., Kawakatsu, H., et
al. (2014). {Changbaishan} volcanism in northeast {China} linked to
subduction-induced mantle upwelling. \emph{Nature Geoscience},
\emph{7}(6), 470--475.

\bibitem[\citeproctext]{ref-tang2023}
Tang, Z., Julià, J., Mooney, W., Mai, P., Yu, H., \& Wu, Y. (2023).
Seismic image of the mantle transition zone beneath northeastern
{China}: Evidence for stagnant {Pacific} subducting slab, lithospheric
delamination and mantle upwelling. \emph{Geophysical Journal
International}, \emph{235}(2), 1872--1887.

\bibitem[\citeproctext]{ref-tauzin2017}
Tauzin, B., Kim, S., \& Kennett, B. (2017). Pervasive seismic
low-velocity zones within stagnant plates in the mantle transition zone:
Thermal or compositional origin? \emph{Earth and Planetary Science
Letters}, \emph{477}, 1--13.

\bibitem[\citeproctext]{ref-vandermeijde2003}
{van der Meijde, M., Marone, F., Giardini, D., \& van der Lee, S.}
(2003). Seismic evidence for water deep in {Earth}'s upper mantle.
\emph{Science}, \emph{300}(5625), 1556--1558.

\bibitem[\citeproctext]{ref-wang2010}
Wang, B., \& Niu, F. (2010). A broad 660 km discontinuity beneath
northeast {China} revealed by dense regional seismic networks in
{China}. \emph{Journal of Geophysical Research: Solid Earth},
\emph{115}(B6).

\bibitem[\citeproctext]{ref-waszek2021}
Waszek, L., Tauzin, B., Schmerr, N., Ballmer, M., \& Afonso, J. (2021).
A poorly mixed mantle transition zone and its thermal state inferred
from seismic waves. \emph{Nature Geoscience}, \emph{14}(12), 949--955.

\bibitem[\citeproctext]{ref-waszek2024}
Waszek, L., Anandawansha, R., Sexton, J., \& Tauzin, B. (2024).
Thermochemistry of the mantle transition zone beneath the western
{Pacific}. \emph{Geophysical Research Letters}, \emph{51}(18),
e2024GL110852.

\bibitem[\citeproctext]{ref-wheeler2014}
Wheeler, J. (2014). Dramatic effects of stress on metamorphic reactions.
\emph{Geology}, \emph{42}(8), 647--650.

\bibitem[\citeproctext]{ref-wheeler2015}
Wheeler, J. (2015). Dramatic effects of stress on metamorphic reactions:
reply. \emph{Geology}, \emph{43}(11), e373--e373.

\bibitem[\citeproctext]{ref-wheeler2020}
Wheeler, J. (2020). A unifying basis for the interplay of stress and
chemical processes in the {Earth}: Support from diverse experiments.
\emph{Contributions to Mineralogy and Petrology}, \emph{175}(12), 116.

\bibitem[\citeproctext]{ref-wu2019}
Wu, W., Ni, S., \& Irving, J. (2019). Inferring {Earth}'s discontinuous
chemical layering from the 660-kilometer boundary topography.
\emph{Science}, \emph{363}(6428), 736--740.

\bibitem[\citeproctext]{ref-xu2008}
Xu, W., Lithgow-Bertelloni, C., Stixrude, L., \& Ritsema, J. (2008). The
effect of bulk composition and temperature on mantle seismic structure.
\emph{Earth and Planetary Science Letters}, \emph{275}(1-2), 70--79.

\bibitem[\citeproctext]{ref-yamazaki2005}
Yamazaki, D., Inoue, T., Okamoto, M., \& Irifune, T. (2005). Grain
growth kinetics of ringwoodite and its implication for rheology of the
subducting slab. \emph{Earth and Planetary Science Letters},
\emph{236}(3-4), 871--881.

\bibitem[\citeproctext]{ref-yu2023}
{Yu, C., Goes, S., Day, E., \& van der Hilst, R.} (2023). Seismic
evidence for global basalt accumulation in the mantle transition zone.
\emph{Science Advances}, \emph{9}(22), eadg0095.

\bibitem[\citeproctext]{ref-zener1946}
Zener, C. (1946). Kinetics of the decomposition of austenite.
\emph{Trans. Aime}, \emph{167}, 550--595.

\end{CSLReferences}


\end{document}
